% ############################################################################
% ##  RETIRED 2026-08-14. NOT COMPILED BY ANYTHING. DO NOT EDIT. DO NOT DELETE.
% ##
% ##  THE THIRD RETIRED MIRROR.  results.tex and assumptions.tex were retired on
% ##  2026-08-03 with a banner of this form; this file is in exactly the same
% ##  position and until now carried only the weaker 2026-08-02 SUPERSEDED note
% ##  below, which was written while the architecture it describes was still
% ##  half-live.  Verified 2026-08-14: NO document \input's this file.  Neither
% ##  paper/qmd/main.qmd nor supplement.qmd references it; scripts/build-theory/
% ##  driver.tex \input's the eight-file split and not this; scripts/build-theory/
% ##  build.sh neither copies nor requires it.  Every mention of the name in the
% ##  repository is a comment, a plan, or a critic report NAMING it as retired.
% ##
% ##  THE 2026-08-02 NOTE BELOW IS STALE IN ONE LOAD-BEARING RESPECT, AND IT IS
% ##  CORRECTED HERE RATHER THAN REWRITTEN, SO THE RECORD STAYS LEGIBLE.  Its
% ##  "LABEL HYGIENE" paragraph states that "results.tex SURVIVES UNDIVIDED and
% ##  is still the source of record: it is what scripts/build-theory/build.sh
% ##  compiles".  THAT HAS BEEN FALSE SINCE 2026-08-03: build.sh compiles the
% ##  SPLIT, results.tex is itself retired, and build.sh:46 records it as such.
% ##  A reader who follows that sentence is sent from one dead file to another.
% ##
% ##  WHY THIS MATTERS MORE THAN "an old copy lying around" -- THE DUPLICATE
% ##  LABELS ARE THE HAZARD, AND THIS FILE HOLDS THE ONES THAT CARRY NUMBERS.
% ##  It declares eq:normconst (:588) and eq:cap (:601), both of which are owned
% ##  by results_s3_estimation.tex (:77, :97) in the compiled corpus.  eq:cap is
% ##  the ONE display in this project whose numerical content moves with the
% ##  reference measure -- see lem:scaleinv(vi)--(vii) and the standing-convention
% ##  paragraph in results_supp.tex -- so a stale second copy of it is the exact
% ##  shape of the two prior divergence incidents: a gate or a reader inspecting
% ##  the wrong artifact and reporting safety, or danger, it never checked.
% ##  It also still declares cor:a1cprim (:352), WITHDRAWN 2026-08-05 for having
% ##  an empty hypothesis set, and still prints the pre-generalisation scalar
% ##  contrast (:1117) and the moment with the contrast on the left (:1120).
% ##
% ##  NOT UPDATED BY THE 2026-08-14 B1(i) PASS, DELIBERATELY.  That pass removed
% ##  the inert "rescaled so that c_Phi = 1" normalisation from B1(i)
% ##  (assumptions_main.tex), added lem:scaleinv(vii) and amended the standing
% ##  convention (results_supp.tex), and retyped three notation_glossary.md rows.
% ##  This file was NOT touched.  It has no c_Phi text to touch -- B1 lives in the
% ##  assumption files -- but its eq:cap and eq:normconst are now formally behind
% ##  the corpus's reading of them.  Behind, not ahead: nothing here is unique
% ##  work that would be lost.
% ##
% ##  KEPT RATHER THAN DELETED on the results.tex / assumptions.tex /
% ##  paper/sections/theorems/_SUPERSEDED.md precedent: nothing in this
% ##  repository is committed, and a deletion would be unrecoverable.
% ############################################################################
%
% SUPERSEDED 2026-08-02.  This file is the FIRST two-way split and is
% referenced by NO .qmd: the article and supplement now input the three-way
% split results_s2_identification.tex / results_s3_estimation.tex /
% results_s4_asymptotics.tex.  Kept only so the split's history is legible.
% Its 8 bare \ref{ass:...} calls were corrected in the same pass as the
% corpus's 82, so build.sh's lint stays green rather than carrying a
% permanent exclusion for a dead file.
%
% ============================================================================
%  RESULTS -- MAIN TEXT.  project `panel-npiv-discrete`
%
%  WHAT THIS FILE HOLDS.  The 23 statements a reader of the article must SEE:
%  the model and the spaces, the operator equation and its linearisation, the
%  identification results, the estimator in all three of its pieces, and the
%  four asymptotic statements.  Nothing else.
%
%  results_main.tex + results_supp.tex ARE the old results.tex.  The split is a
%  PARTITION: all 85 theorem-like environments survive, every \label is declared
%  exactly once across the two files, and no statement was rewritten, retitled or
%  renumbered.  Each environment carries the comment block that preceded it.
%
%  results_supp.tex holds the other 62 environments, the original file header,
%  and the change log.  It is \input by paper/qmd/supplement.qmd ahead of
%  proofs.tex; this file is \input by paper/qmd/main.qmd.
%
%  LABEL HYGIENE -- READ THIS BEFORE ADDING AN \input ANYWHERE.
%  results.tex SURVIVES UNDIVIDED and is still the source of record: it is what
%  scripts/build-theory/build.sh compiles, what scripts/render-sync/
%  check_rendering_sync.py reads, and what ~100 files under scripts/symbolic/
%  specs/ cite BY LINE NUMBER.  It was therefore not edited.  The consequence is
%  that every \label in this file is declared TWICE on disk, here and in
%  results.tex.  THIS FILE AND results.tex MUST NEVER BE \input INTO THE SAME
%  DOCUMENT, and neither must results_supp.tex and results.tex.  That is the same
%  rule results.tex's own header states for paper/sections/theorems/*.tex.
%
%  CROSS-FILE REFERENCES.  A statement kept here may \cref a statement that moved
%  to the supplement, and vice versa.  Those calls do not resolve, because the
%  article and the supplement are separate documents.  The MANUSCRIPT PROSE never
%  does this -- every \cref and \eqref in main.qmd points at a label defined in
%  this file or in assumptions.tex.  The residue is internal to the statements,
%  which are preserved verbatim by instruction.
% ============================================================================



% ===========================================================================
%  PART A.  SETUP
% ===========================================================================

\begin{definition}[Model and differencing]\label{def:model}
The structural model is
\begin{equation}\label{eq:levels}
Y_{ij} \;=\; \vdagger(X_{ij}) \;+\; Z_{ij}'\theta_\dagger \;+\; \eta_i \;+\; U_{ij},
\qquad j\in\{1,2\},
\end{equation}
with $X_{ij}\in\mathcal X\subset\IR$ scalar and continuously distributed, $Z_{ij}\in\IR^{p}$,
$W_i\in\{0,1\}$, and $\eta_i$ an individual effect that may depend on $(X,Z,W)$. Differencing across
the two periods,
\begin{equation}\label{eq:diffs}
% \Delta\vdagger_i was a DOUBLE SUBSCRIPT error: \vdagger expands to
% \varphi_{\dagger}, which already carries one.  Written out here; the reading is
% fixed by the convention on the next line.  See preamble.tex at \vdagger.
\Delta Y_i \;=\; \Delta\varphi_{\dagger i} \;+\; \Delta Z_i'\theta_\dagger \;+\; \Delta U_i,
\qquad \Delta\varphi_i:=\varphi(X_{i2})-\varphi(X_{i1}),
\end{equation}
and, for a trial pair $(\varphi,\theta)$, the residual at trial values is
\begin{equation}\label{eq:trialres}
\Delta U_i(\theta,\varphi)\;:=\;\Delta Y_i-\Delta\varphi_i-\Delta Z_i'\theta ,
\qquad \Delta U_i \;=\; \Delta U_i(\theta_\dagger,\vdagger).
\end{equation}
The subscript $i$ is suppressed whenever a statement concerns the population distribution $P$.
\end{definition}

\begin{definition}[Spaces, and the domain of the operator]\label{def:spaces}
Let $\pi$ be the reference density on $\mathcal X$ of Assumption~C1, and set
\[
\mathcal V\;:=\;\Big\{\varphi\in L^{2}(\mathcal X,\pi):\ \textstyle\int_{\mathcal X}\varphi\,d\pi=0\Big\},
\qquad
L^{2}(\Pi)\;:=\;L^{2}\big(\IR\times\mathcal T,\ \Pi\big),
\]
with $\Pi$ as in Assumption~B1(iv). Inner products are
$\langle h,h'\rangle_{L^2(\pi)}=\int hh'\,\pi\,dx$ and
$\langle\Xi,\Xi'\rangle_{L^2(\Pi)}=\int\Xi\Xi'\,d\Pi$ --- \emph{the bound variables are $h,h'$ on the
domain side and $\Xi,\Xi'$ on the range side, being this project's letters for a direction in
$L^{2}(\pi)$ and a generic element of a range-type space respectively}. The period-difference
% ARCHAEOLOGY (moved out of compiled text, shortening pass 2026-08-01).
% A round-11 draft wrote the two pairings as <a,b>, with a the RESERVED degree of
% ill-posedness and b the weight anchor b_star -- glossary C72's species, in the
% one display that defines the pairing.
operator is
$\Delta:\mathcal V\to L^{2}(P)$, $\Delta h=h(X_2)-h(X_1)$. For an integrable random variable $g$ write
\begin{equation}\label{eq:nudef}
\nu_g(x)\;:=\;E[g\mid X_2=x]\,f_{X_2}(x)\;-\;E[g\mid X_1=x]\,f_{X_1}(x),
\qquad
\nu\;:=\;\nu_{\Delta Z}\in\IR^{p} ,
\end{equation}
the latter understood coordinatewise.

\emph{Every operator in this paper has domain $\mathcal V$, with one declared exception named below.} The
orthogonal decomposition $L^{2}(\pi)=\mathcal V\oplus\{\text{constants}\}$ holds because
$\langle c,h\rangle_{L^2(\pi)}=c\int h\,d\pi$, and $\Delta$ annihilates constants, so nothing is lost by
the restriction while a zero singular value is avoided. In particular the singular system
$(\sigma_j;e_j,\tilde e_j)_{j\ge1}$ of \cref{prp:compact}, the completeness condition A1c, the source
condition A2 and the Hilbert-scale link H1 all refer to
$\mathcal A_{\mathrm{cond}}:\mathcal V\to L^{2}(\Pi)$ and to that same singular system.

\emph{The exception is the second-order form operator $\mathcal B$ of \eqref{eq:Bpi}, which is read on
all of $L^{2}(\pi)$}, because $\mathcal V$ is \emph{not} invariant under multiplication by a non-constant
function and $\mathcal B$ contains such a multiplication. The quadratic form of
\cref{lem:secondord}(i) is evaluated only at $h\in\mathcal V$, so the form's own sharp constant is
governed by the compression $P_{\mathcal V}\mathcal B_{\mathrm{sym}}P_{\mathcal V}|_{\mathcal V}$, where
$P_{\mathcal V}$ is the $L^{2}(\pi)$-orthogonal projection onto $\mathcal V$; \cref{lem:solset}(iii)
records the relation between the two. \emph{No other object in the paper leaves $\mathcal V$.}
\end{definition}

% ---------------------------------------------------------------------------
%  THE OPERATOR EQUATION.  This is the object missing from the source draft at
%  `:120--122', where the sentence ends "Then we have that" and the display
%  never appears.  Everything in draft Section 2 after that point -- the
%  Frechet derivative at `:124--126', the completeness condition at
%  `:127--130', the adjoint at `:172--177' -- is a derivative of this
%  definition and cannot be read without it.
% ---------------------------------------------------------------------------

\begin{definition}[The operator equation]\label{def:operator}
Let $\Delta Y_\theta:=\Delta Y-\Delta Z'\theta$, so that $\Delta U(\theta,\varphi)=\Delta Y_\theta-\Delta\varphi$.
For $(\varphi,\theta)\in\mathcal V\times\IR^{p}$ define the \emph{Bierens moment family}
\begin{equation}\label{eq:Mdef}
\mathcal M(u,t;\varphi,\theta)
\;:=\;E\Big[\widetilde W\;\un\big\{\Delta Y_\theta-\Delta\varphi\le u\big\}\;w\big(t'\Phi(\Delta Z)\big)\Big],
\qquad (u,t)\in\IR\times\mathcal T,
\end{equation}
with $\widetilde W=W-\varrho(\Delta Z)$, and let
\begin{equation}\label{eq:Tdef}
T:\ \mathcal V\times\IR^{p}\longrightarrow L^{2}(\Pi),
\qquad
T(\varphi,\theta)\;:=\;\big[(u,t)\mapsto\mathcal M(u,t;\varphi,\theta)\big].
\end{equation}
The \emph{operator equation of the model} is
\begin{equation}\label{eq:opeq}
T(\varphi,\theta)\;=\;0 \quad\text{in }L^{2}(\Pi).
\end{equation}
We abbreviate $\mathcal M(u,t):=\mathcal M(u,t;\vdagger,\theta_\dagger)$ and
$T(\varphi):=T(\varphi,\theta_\dagger)$ when no ambiguity arises.
\end{definition}

\begin{definition}[Admissible sets]\label{def:admissible}
Restrictions (A) and (B) of plan \S2.1 are
\[
\text{(A)}\quad E\big[\widetilde W\,\varsigma(\Delta U)\;\big|\;\Delta Z\big]=0\ \ \text{a.s. for every bounded measurable }\varsigma,
\qquad
\text{(B)}\quad E\big[\Delta Z\,\Delta U\big]=0 .
\]
Define
\[
\mathcal I\;:=\;\Big\{(\varphi,\theta)\in\mathcal V\times\IR^{p}:\ \Delta Y-\Delta\varphi-\Delta Z'\theta\ \upmodels\ W\ \big|\ \Delta Z\Big\},
\qquad
\mathcal V_{CI}\;:=\;\big\{\varphi\in\mathcal V:\ (\varphi,0)\in\mathcal I\big\}.
\]
Conditional independence $\xi\upmodels W\mid\Delta Z$ is understood in the factorisation sense,
$P(\xi\in B,\,W=j\mid\Delta Z)=P(\xi\in B\mid\Delta Z)\,P(W=j\mid\Delta Z)$ a.s.\ for all Borel $B$ and
$j\in\{0,1\}$. This is equivalent to the conditional-expectation form when both $W$-cells are non-null
given $\Delta Z$, and it is the form that behaves correctly when $W$ is conditionally degenerate --- in
which case it holds automatically, a fact used in the proof of \cref{lem:bierens}.
\end{definition}

\begin{proposition}[Fr\'echet derivative in $\varphi$, with a Lipschitz derivative map]\label{prp:frechet}
$T$ is Fr\'echet differentiable in $\varphi$ with a quadratic remainder, and the derivative map is
Lipschitz.
Suppose B1(i), (ii) and (iv.c), C1 and C2$'$(i) hold. Then, for every $\theta\in\IR^{p}$ and every
$\varphi\in\mathcal V$, the linear operator $\mathcal A_\varphi:\mathcal V\to L^{2}(\Pi)$ defined by
\begin{equation}\label{eq:Adef}
\mathcal A_\varphi h(u,t)
\;=\;E\Big[\widetilde W\,w\big(t'\Phi(\Delta Z)\big)\,
f_{\Delta Y_\theta-\Delta\varphi\mid X_1,X_2,\Delta Z,W}(u)\;\big\{h(X_2)-h(X_1)\big\}\Big]
\end{equation}
is bounded, with $\|\mathcal A_\varphi\|\le2\,\bar c_\pi^{1/2}\|\bar f_0\|_\infty c_\Pi$, and:
\begin{enumerate}
\item[(i)] \emph{Fr\'echet differentiability with a quadratic remainder:} for all $h\in\mathcal V$,
\begin{equation}\label{eq:frechetbound}
\big\|T(\varphi+h,\theta)-T(\varphi,\theta)-\mathcal A_\varphi h\big\|_{L^{2}(\Pi)}
\;\le\;C_{\!\mathcal A}\,\|h\|_{L^{2}(\pi)}^{2},
\qquad C_{\!\mathcal A}:=2\,\bar c_\pi\,\|\bar f_1\|_\infty\,c_\Pi ;
\end{equation}
\item[(ii)] \emph{Lipschitz derivative map:} for all $h\in\mathcal V$,
\begin{equation}\label{eq:lipschitz}
\big\|\mathcal A_{\varphi+h}-\mathcal A_\varphi\big\|\;\le\;2\,C_{\!\mathcal A}\,\|h\|_{L^{2}(\pi)} .
\end{equation}
\end{enumerate}
At the truth we write $\mathcal A_{\mathrm{cond}}:=\mathcal A_{\vdagger}$ with $\theta=\theta_\dagger$,
in which case the conditional density in \eqref{eq:Adef} is $f_{\Delta U\mid X_1,X_2,\Delta Z,W}$.
\end{proposition}

\begin{proposition}[Compactness]\label{prp:compact}
$\mathcal A_{\mathrm{cond}}$ is Hilbert--Schmidt, hence compact.
Suppose B1(i), (ii) and (iv.c), C1 and C2$'$(i) hold. Then $\mathcal A_{\mathrm{cond}}$ is the integral
operator on $\mathcal V$ with kernel
\begin{equation}\label{eq:kernel}
\mathsf k(u,t;x)\;=\;\frac{\nu_{g_{u,t}}(x)}{\pi(x)},
\qquad
g_{u,t}\;:=\;\widetilde W\,w\big(t'\Phi(\Delta Z)\big)\,f_{\Delta U\mid X_1,X_2,\Delta Z,W}(u),
\end{equation}
and $\|\mathcal A_{\mathrm{cond}}\|_{\mathrm{HS}}\le 2\,\bar c_\pi\,\|\bar f_0\|_\infty\,c_\Pi<\infty$.
Moreover $\mathsf k(u,t;\cdot)\in\mathcal V$ for every $(u,t)$, so $\mathcal A^\ast_{\mathrm{cond}}$ maps
$L^{2}(\Pi)$ into $\mathcal V$; and $\mathcal A_{\mathrm{cond}}$ admits a singular system
$(\sigma_j;e_j,\tilde e_j)_{j\in\mathcal J}$, $\mathcal J\subseteq\IN$, with $\{e_j\}_{j\in\mathcal J}$
an orthonormal basis of $(\ker\mathcal A_{\mathrm{cond}})^{\perp}\cap\mathcal V$ and, when $\mathcal J$ is
infinite, $\sigma_j\downarrow0$.
\end{proposition}

% prp:compact's own hypotheses do NOT include A1b(iv): compactness, the kernel
% and the singular system need nothing about Delta Z's moments.  The nu-versus-
% nu/pi comparison does need E||Delta Z|| < infinity, so it is lifted out of the
% proof into rem:nupi, where the assumption can be cited.  Round 2 ran the
% comparison inside prp:compact's proof without listing it.
%
% THE RANK CLAUSE IS NEW.  Both counterexample-hunter DGPs produce a RANK-ONE
% A_cond satisfying every hypothesis above, so finite rank is a live case, not
% a pathology.  Nothing here breaks -- but the singular system then spans
% (ker A_cond)^perp, NOT V, and H1's sigma_j asymp j^{-a} is false, which is why
% H1 now assumes infinite rank explicitly.

% prp:compact closes the gap the round-3 strategist-critic recorded at m-23:
% compactness of A_cond is a PREMISE of prp:sharpquad(ii) and of every Picard,
% source-condition and Landweber statement in the project, and it appeared
% nowhere.  It is proved here rather than assumed.
%
% ROUND 2: the clause "k(u,t;.) in V" is new and is not decoration.  It is the
% internal check that A2 is stated on the right object.  Every element of
% R(A*_cond) is an integral of kernel slices; every slice lies in V; hence
% R(A*_cond) is contained in V, which is exactly {zeta : int zeta pi dx = 0}.
% And nu_g/pi lies in V for every g (int (nu_g/pi) d pi = int nu_g dx = 0),
% while nu_g does not.  See the note under A2 in assumptions.tex.
%
% NOTE THE REASON, which changed in the cleanup pass.  The containment
% R(A*) subset V follows from the CONSTRUCTION of the adjoint, not from
% R(A*) subset (ker A)^perp together with Delta 1 = 0: once def:spaces fixes
% the domain at V, the constants are not in the domain at all and that
% orthogonality argument is unavailable.  The conclusion is unchanged.

% ===========================================================================
%  PART B.  STEP 1 -- SHIFT INVARIANCE
% ===========================================================================

\begin{lemma}[L1: shift invariance]\label{lem:shiftinv}
Restriction (A) is invariant to the value of $\theta$, so it does not identify $\theta_\dagger$.
The set $\mathcal I$ of \cref{def:admissible} satisfies, with no assumption beyond that definition,
\begin{equation}\label{eq:shiftinv}
\mathcal I\;=\;\mathcal V_{CI}\times\IR^{p};
\end{equation}
and if in addition A1 holds, then $\vdagger\in\mathcal V_{CI}$, so $(\vdagger,\theta)\in\mathcal I$ for
every $\theta\in\IR^{p}$.
\end{lemma}

% A4(i) IS NOT A HYPOTHESIS HERE, and its removal is a cleanup-pass correction.
% \eqref{eq:binaryquant} holds for ANY version of the conditional distribution
% functions, including on the degenerate region D := {z : varrho(z)(1-varrho(z))
% = 0} where one of the two W-cells is null and F_j(.|z) is an arbitrary
% version: on D the prefactor varrho(1-varrho) is zero, so both sides vanish
% whatever version is used.  The extension step needs only the POINTWISE
% positivity varrho(z)(1-varrho(z)) > 0, which the proof states in exactly that
% form, and it draws no conclusion on D -- where none is needed, since W is
% conditionally degenerate there and lem:bierens handles it separately.

\begin{lemma}[Bierens equivalence for the panel restriction]\label{lem:bierens}
The $\Delta Z$-conditional restriction (A) is equivalent to the vanishing of the unconditional family
$\{\mathcal M(u,t)\}$, so no conditional distribution given $\Delta Z$ ever has to be estimated.
Suppose B1(i)--(iv) hold. Then the following are equivalent:
\begin{enumerate}
\item[(a)] $\Delta U\upmodels W\mid\Delta Z$;
\item[(b)] $E[\widetilde W\,\varsigma(\Delta U)\mid\Delta Z]=0$ a.s.\ for every bounded measurable $\varsigma$, that is, restriction \textup{(A)};
\item[(c)] $\mathcal M(u,t)=0$ for every $(u,t)\in\IR\times\mathcal T$;
\item[(d)] $T(\vdagger,\theta_\dagger)=0$ in $L^{2}(\Pi)$.
\end{enumerate}
\end{lemma}

\begin{corollary}[to \cref{lem:shiftinv} and \cref{lem:bierens}]\label{cor:opsol}
The operator equation has the same solution set in $\varphi$ at every $\theta$, and that set is
$\mathcal V_{CI}$.
Suppose B1(i)--(iv) hold. Then for every $\theta\in\IR^{p}$,
\[
\big\{\varphi\in\mathcal V:\ T(\varphi,\theta)=0\ \text{in }L^{2}(\Pi)\big\}\;=\;\mathcal V_{CI},
\]
and under A1b(iii) this set equals $\{\vdagger\}$.
\end{corollary}

\begin{proposition}[R2: when a weight yields a $\varphi$-free moment]\label{prp:nophifree}
A moment is free of the structural function exactly when the two periods' $g$-tilted marginals of the
endogenous regressor coincide.
Suppose C1(ii) holds and let $g$ be a random variable satisfying
\begin{equation}\label{eq:gproviso}
E|g|<\infty,
\qquad
E\big|g\,h(X_j)\big|<\infty\ \ (j=1,2)\ \text{for every }h\in\mathcal V,
\qquad
E\big|g\,\Delta U(\theta,0)\big|<\infty .
\end{equation}
Then the map $\varphi\mapsto E\big[g\,\Delta U(\theta,\varphi)\big]$ is constant on $\mathcal V$, for one
and hence every $\theta\in\IR^{p}$, if and only if
\begin{equation}\label{eq:tilted}
E[g\mid X_2=x]\,f_{X_2}(x)\;=\;E[g\mid X_1=x]\,f_{X_1}(x)
\qquad\text{for almost every }x\in\mathcal X,
\end{equation}
that is, if and only if $\nu_g\equiv0$.

% ARCHAEOLOGY.  A1b(i) was listed as a hypothesis here through round 7 and is struck.
\emph{A1b(i) is \textbf{not} a hypothesis.} The proof's third step upgrades
``$\nu_g/\pi\perp\mathcal V$'' to ``$\nu_g\equiv0$'' using $L^{2}(\pi)=\mathcal V\oplus\{\text{constants}\}$
--- which is \cref{def:spaces} --- together with $\int_{\mathcal X}\nu_g\,dx=0$ from
\cref{lem:nusigma}(iv). $\vdagger$ occurs nowhere in the statement or the proof, and the map is
quantified over generic $\varphi,\varphi_0\in\mathcal V$.
\end{proposition}

\begin{corollary}[to \cref{prp:nophifree}: the debiasing weight must vary with the residual]\label{cor:udep}
If the debiasing weight does not depend on the residual, the debiased moment is not merely
first-order insensitive to the structural function but exactly free of it.
Suppose A1, A1b(i), A1b(iv) and C1(ii) hold, let $L$ be differentiable in $u$ \emph{at every} $(u,z)$
with $\partial_uL(u,z)=c(z)$ for some $c$ such that $g:=\Delta Z+\widetilde W\,c(\Delta Z)$ satisfies
\eqref{eq:gproviso} coordinatewise, and suppose $L$ satisfies N1. Then
$L(u,z)=c(z)\,u$, the moment function is $\psi=g\,\Delta U(\theta,\varphi)$ with $g$ measurable with
respect to $\sigma(W,\Delta Z)$, and the following are equivalent:
\begin{enumerate}
\item[(a)] $E[\psi]$ is first-order insensitive to $\varphi$ at $\vdagger$, i.e.\
$\nu_{\Delta Z+\widetilde Wc}\equiv0$;
\item[(b)] $\varphi\mapsto E[\psi]$ is constant on $\mathcal V$;
\item[(c)] $\nu_{\widetilde Wc}=-\nu$ \emph{for the $c$ of the hypothesis}; equivalently, membership
$-\nu\in\mathcal N_0:=\{\nu_{\widetilde W\,\tilde c(\Delta Z)}:\tilde c\in\mathcal C_0\}$ is realised
\emph{at that} $c$, where $\mathcal C_0$ is the set of $\tilde c$ for which
$\Delta Z+\widetilde W\tilde c(\Delta Z)$ satisfies \eqref{eq:gproviso} coordinatewise.
\end{enumerate}
Moreover $\varphi\mapsto E[\psi]$ is affine on $\mathcal V$, so every term of order two or higher in its
expansion vanishes identically \emph{whether or not \textup{(a)--(c)} hold}; what \textup{(a)--(c)} add is
that the first-order term vanishes as well. Contrapositively, if $-\nu\notin\mathcal N_0$, then every
$\ell_\star$ satisfying $\nu_{\Delta Z+\widetilde W\ell_\star}\equiv0$ has
$\partial_u\ell_\star\not\equiv0$.
\end{corollary}

\begin{proposition}[Characterisation of the kernel of the linearisation]\label{prp:kerchar}
A direction is annihilated by the linearised operator exactly when the instrument fails to move its
period difference, given the residual and the controls.
Suppose B1(i)--(iv), C1 and C2$'$(i) hold, and let $h\in\mathcal V$. Then
\begin{enumerate}
\item[(i)] the following are equivalent:
\emph{(a)} $\mathcal A_{\mathrm{cond}}h=0$ in $L^{2}(\Pi)$;
\emph{(b)} $E\big[\widetilde W\,\{h(X_2)-h(X_1)\}\,\big|\,\Delta U,\Delta Z\big]=0$ almost surely;
\item[(ii)] if in addition A1 and A4(i) hold, \emph{(b)} is equivalent to
\emph{(c)} $E\big[h(X_2)-h(X_1)\mid\Delta U,\Delta Z,W=1\big]
=E\big[h(X_2)-h(X_1)\mid\Delta U,\Delta Z,W=0\big]$ almost surely.
\end{enumerate}
\end{proposition}

\begin{corollary}[to \cref{prp:kerchar}(i): a primitive for A1c]\label{cor:a1cprim}
Instrument completeness given the residual, together with measure separability, delivers first-order
local identification.
Suppose B1(i)--(iv), C1, C2$'$(i), A1b(ii) and A1d hold. Then A1c holds, that is,
$\ker\mathcal A_{\mathrm{cond}}=\{0\}$. \emph{A1b(i) is \textbf{not} a hypothesis}: the ``constant,
hence zero'' step of \cref{lem:nusigma}\textup{(ii)} is \cref{def:spaces} and not that clause, and the
conclusion is about a generic $h\in\ker\mathcal A_{\mathrm{cond}}\subseteq\mathcal V$.
% ARCHAEOLOGY.  A1b(i) was listed here through round 7, inherited from
% lem:nusigma(ii), and is struck.
\end{corollary}

% ===========================================================================
%  PART F.  STEP 5 -- THE SELECTION, ORTHOGONALITY, AND THE BALANCING BOUND
%
%    def:weight    representers, weights, and THE ANCHOR (round-7: closes the
%                  CRITICAL that L_star was used everywhere and defined nowhere)
%    def:dirclass  admissible perturbation directions, with the SHIFT-UNIFORM
%                  seminorms that make m(tau) exist off tau_1 = 0
%    lem:solset   (L-16)  structure of the solution set and the selection
%    lem:secondord (L-6)  the second-order expansion in three nuisances
%    prp:orthog           orthogonality in varphi, in L, and in varrho
%    lem:balance   (L-7)  the balancing bound on the first-order bias
%
%  TYPING.  Per F-15 every statement below is COORDINATEWISE in the p
%  coordinates of theta: L, ell, ell_star and the directions D are R^p-valued,
%  and A*_cond acts on each coordinate separately.  eq:JSig's J is the one
%  object whose FORM is NOT coordinatewise; see the round-7 note in
%  assumptions.tex.  DISPATCH 13: the tag on that sentence was "A3", which is
%  now retired (ass:rank), and the distinction between FORM and VALUE is now
%  the operative one -- lem:momzero(iv) makes the non-coordinatewise summand
%  ZERO, so J itself IS coordinatewise and symmetric, while J-hat (eq:Jhat) is
%  neither.  The typing warning therefore transfers from J to J-hat and is not
%  withdrawn.
%
%  SYMBOLS INTRODUCED OR MOVED IN ROUND 7, all recorded in the glossary:
%    Psi(tau)  replaces m(tau).  C7 retired the shorthand m := M/pi BECAUSE it
%              collided with the LF iteration counter m, and round 6 introduced a
%              second collision on the same letter.  Psi is the expectation of
%              psi, which is the mnemonic; the pair is documented in Sec.1.6.
%    sans-r    replaces r as the varrho-direction.  Italic r is the propensity
%              smoothness index and the two CO-OCCUR in one display of
%              lem:secondord(iv)'s proof, where r_n = n^{-r/(2r+d_Z)}.
%    tau       is Part F's perturbation parameter; A7's discrepancy constant has
%              moved to tau_DP.
%    vartheta  is the R^p coordinate-contraction direction, replacing the
%              round-6 binding use of varsigma, which is a RESERVED LOCAL bound
%              variable and could not be made binding.
%    D         RETAINED, and C2's retirement of plain D is LIFTED -- see the
%              glossary Sec.8.6.  Three fonts, three types: calligraphic D the
%              differentiation operator, bold D the sieve design matrix, italic D
%              the direction in L.  This is the rule Sec.8.2 already applies to
%              B_sym versus bold B.
% ===========================================================================

\begin{definition}[Representers, weights, and the anchor]\label{def:weight}
Let $\ell^\Pi\in L^{2}(\Pi)$ and $\ell:=\mathcal W\ell^\Pi$ with $\mathcal W$ as in
\eqref{eq:Wdef}, applied coordinatewise.
\begin{enumerate}
\item[(a)] $\ell^\Pi$ is an \emph{admissible representer} if it satisfies \textbf{A2b}. \emph{That tag,
and not A$2'$(i), is what this clause names}: A$2'$(i) is a \emph{maintained assertion} about
$\ell^\Pi_\star$; A2b is a \emph{property} a candidate representer may or may not have, and it is the
property every result quantifying over representers imposes. Under B1(i), (ii), (iv.a),
\cref{cor:elldiff} --- itself stated at a generic admissible representer --- applies to $\ell$: it is
bounded and continuously differentiable in $u$, uniformly in $z$, with budgets
$\Lambda^{(0)},\Lambda^{(1)}$.
\item[(b)] For a Borel $b:\IR^{p}\to\IR^{p}$, the \emph{weight generated by $\ell$ and anchored at $b$}
is
\begin{equation}\label{eq:weightdef}
L^{\,b}_{\ell}(u,z)\;:=\;\int_{0}^{u}\ell(v,z)\,dv\;+\;b(z),
\qquad (u,z)\in\IR\times\IR^{p},
\end{equation}
the integral being well defined because $\ell(\cdot,z)$ is continuous and bounded. It is the unique
function that is differentiable in $u$ at \emph{every} point with $\partial_uL^{\,b}_\ell=\ell$ and
$L^{\,b}_{\ell}(0,\cdot)=b$. The weight is \emph{admissible} if $E\|b(\Delta Z)\|^{2}<\infty$.
\item[(c)] $\ell^\Pi$ is an \emph{admissible selection} if it is an admissible representer and solves
$\mathcal A^\ast_{\mathrm{cond}}\ell^\Pi=-\nu/\pi$.
\end{enumerate}
Under A2 and A2$'$(i), $\ell^\Pi_\star$ is an admissible selection --- A2 gives that it solves the
adjoint equation, A2$'$(i) that it satisfies A2b --- with $\ell_\star=\mathcal W\ell^\Pi_\star$, and
$L_\star:=L^{\,b_\star}_{\ell_\star}$ is the admissible weight fixed by A2$'$(iii). Write
$\zeta_\star:=(\vdagger,L_\star,\varrho)$ for the nuisance triple at the truth.
\end{definition}

\begin{proposition}[Neyman orthogonality]\label{prp:orthog}
The debiased moment is orthogonal in all three nuisances: in the structural function exactly when the
debiasing weight solves the adjoint equation, and in the weight and the propensity exactly, for every
admissible direction.
Suppose A1b(i), A1b(iv), B1(i), (ii), (iv), C1(ii) and C2$'$(i) hold. \emph{A1 is \textbf{not} in the
global list}: it is consumed by part (ii) alone --- where it is
assumed and cited --- while parts (i) and (iii) use it nowhere. The split follows the precedent already
in this file at \cref{lem:shiftinv}, and its practical point is the same: \cref{lem:balance}, which
inherits ``the hypotheses of \cref{prp:orthog}(i)'', would otherwise inherit A1 as surplus and consume
it in no step.
\emph{N1 is not assumed in any part}: (i) and (ii) do not use it and (iii) characterises it. Then
\begin{enumerate}
\item[(i)] \emph{Structural function --- stated at an arbitrary admissible representer, and A2 is not
assumed.} Let $\ell^\Pi$ be an admissible representer in the sense of A2b,
$\ell:=\mathcal W\ell^\Pi$, and $L$ an admissible
weight generated by $\ell$ in the sense of \cref{def:weight}. For $h\in\mathcal V$ define
\begin{equation}\label{eq:PsiL}
\Psi_L(\tau_1)\;:=\;E\big[\psi\big(\theta_\dagger,(\vdagger+\tau_1h,\,L,\,\varrho)\big)\big]
\;=\;E\Big[\Delta Z\,(\Delta U-\tau_1\Delta h)+\widetilde W\,L\big(\Delta U-\tau_1\Delta h,\Delta Z\big)\Big],
\qquad\tau_1\in\IR,
\end{equation}
which is finite for \emph{every} $\tau_1\in\IR$: the second term because
$\|L(u,z)\|\le\|b(z)\|+\|\ell\|_\infty|u|$ by \eqref{eq:weightdef}, with $|\widetilde W|\le1$ by (P1),
$E\|b(\Delta Z)\|^{2}<\infty$ by admissibility, $\|\ell\|_\infty<\infty$ by \cref{cor:elldiff} under A2b
--- \emph{the tag is A2b and not A$2'$(i), because $\ell^\Pi$ here is a generic admissible representer
and not $\ell^\Pi_\star$} --- and
$E|\Delta U|,E|\Delta h|<\infty$ by A1b(iv) and \cref{lem:nusigma}(i); the first because
$E\|\Delta Z\Delta U\|<\infty$ and
$E\|\Delta Z\Delta h\|\le\|\Delta Z\|_{L^{2}(P)}\|\Delta h\|_{L^{2}(P)}<\infty$ by A1b(iv), the
Cauchy--Schwarz inequality and \cref{lem:nusigma}(i). \emph{This defines $\Psi_L$ here
and does not import it}: no radius $\bar\tau$, no direction pair $(D,\mathsf r)$ and no clause of
A2$'$ enters \eqref{eq:PsiL}, which is what makes the hypothesis reduction of this clause readable off
its own statement. At $L=L_\star$, $\Psi_{L_\star}(\tau_1)=\Psi(\tau_1,0,0)$ in the notation of
\cref{lem:secondord}, so the two agree wherever both are defined. Then for every $h\in\mathcal V$,
$\partial_{\tau_1}\Psi_L(0)=-\big\langle\nu_{\Delta Z+\widetilde W\ell}/\pi,\,h\big\rangle_{L^{2}(\pi)}$,
and this vanishes for every $h\in\mathcal V$ if and only if
\begin{equation}\label{eq:orthphi}
\nu_{\Delta Z+\widetilde W\ell}\equiv0
\qquad\Longleftrightarrow\qquad
\mathcal A^\ast_{\mathrm{cond}}\ell^\Pi=-\nu/\pi .
\end{equation}
\emph{In particular, under A2$'$ the left-hand condition holds at $\ell=\ell_\star$}, and
$\partial_{\tau_1}\Psi(0)=0$ for every $h\in\mathcal V$.
\item[(ii)] \emph{Debiasing weight, exactly.} Assume in addition A1, A2, A2$'$(i) and A2$'$(iii).
\emph{This is the one part that consumes A1}, through the conditional independence that makes
$E[\widetilde W\,D\mid\Delta Z]=E[\widetilde W\mid\Delta Z]\,E[D\mid\Delta Z]=0$. For every $D$ with
$E[\|D\|\mid\Delta Z]<\infty$ almost surely and $E\|\widetilde WD\|<\infty$, and every $\tau_2\in\IR$,
$E[\psi(\theta_\dagger,(\vdagger,L_\star+\tau_2D,\varrho))]=E[\psi(\theta_\dagger,\zeta_\star)]$.
\item[(iii)] \emph{Propensity: exact orthogonality is equivalent to N1.} Assume in addition A2,
A2$'$(i), A2$'$(iii) and A4(i), and
\emph{do not assume} N1 --- nor A1, which this part does not use. Then, for every bounded $\mathsf r$
with $\varrho+\tau_3\mathsf r$ taking values
in $[0,1]$ for all small $\tau_3$, and every such $\tau_3$,
\[
E\big[\psi(\theta_\dagger,(\vdagger,L_\star,\varrho+\tau_3\mathsf r))\big]=E\big[\psi(\theta_\dagger,\zeta_\star)\big]
\qquad\Longleftrightarrow\qquad
E\big[L_\star(\Delta U,\Delta Z)\mid\Delta Z\big]=0\ \text{a.s.},
\]
the right-hand condition being exactly N1 --- \emph{and both sides are well defined, because A2$'$(iii)
with \cref{cor:elldiff} and A1b(iv) give $E[\|L_\star\|\mid\Delta Z]<\infty$ almost surely.}
\end{enumerate}
\end{proposition}

\begin{proposition}[The second-order remainder on the ball: the sharp constant, and the exact $\delta^{2}$ order under O-16(a)]\label{prp:sharpquad}
The supremum of the second-order remainder over a ball of radius $\delta$ in the structural direction is
exactly $\tfrac12\delta^{2}$ times the norm of the compressed symmetrised form, and it is bounded below by
$\tfrac12\delta^{2}\|\vartheta'\sigma_q/\pi\|_\infty$. \emph{So, provided the form does not degenerate ---
which is a hypothesis and is stated as one} --- no bound of the form $o(\delta^{2})$, and in particular no
bound in the weak norm $\|\mathcal A_{\mathrm{cond}}\cdot\|_{L^{2}(\Pi)}$, holds uniformly on that ball.
Suppose the hypotheses of \cref{lem:solset} hold --- \emph{cited by clause and not as a block}: clause
\textup{(iv)} in addition requires \textbf{C2$'$(i)}, which it consumes transitively through
\cref{lem:formstruct}\textup{(ii)} and \cref{prp:compact}, for the compactness of
$\mathcal A_{\mathrm{cond}}$ that delivers $\|\mathcal A_{\mathrm{cond}}h_k\|\to0$; \emph{clauses
\textup{(i)}--\textup{(iii)} need it not}, and Step 5 of the proof runs the implication
O-16\textup{(b)}$\Rightarrow$O-16\textup{(a)} through \eqref{eq:complower} and the numerical radius
precisely so that they need not. Fix $\vartheta\in\IR^{p}$ with $\|\vartheta\|=1$,
and let the \emph{second-order remainder in the structural direction} at $h\in\mathcal V$ be
\begin{equation}\label{eq:quadrem}
\tfrac12\,\vartheta'\partial^{2}_{\tau_1\tau_1}\Psi(0)
\;=\;\tfrac12\,E\big[\widetilde W\,\vartheta'\partial_u\ell_\star\,(\Delta h)^{2}\big]
\;=\;\tfrac12\,\big\langle\mathcal Bh,h\big\rangle_{L^{2}(\pi)} ,
\end{equation}
which is \cref{lem:secondord}\textup{(i)} contracted in $\vartheta$ and transferred to $L^{2}(\pi)$ by
\cref{lem:solset}\textup{(iii)}. \emph{\Cref{lem:secondord} is read at the admissible direction pair
$(D,\mathsf r)=(0,0)$}, which is admissible at every $(h,\bar\tau)$ by \cref{def:dirclass} and is
differentiable there, so that $\Psi(\cdot,0,0)=\Psi_{L_\star}$ by \eqref{eq:PsiL} and \eqref{eq:quadrem} is
half the second $\tau_1$-derivative of the moment along the \emph{structural direction alone}; no
condition on a direction in $L$ or in $\varrho$ is inherited. \emph{The remainder carries no symbol of
its own}; italic $Q$ names the regular conditional distributions of \textup{(P6b)}.
Then, for every $\delta>0$:
\begin{enumerate}
\item[(i)] \emph{The sharp constant, with no hypothesis on $q$ beyond \cref{lem:solset}'s.}
\begin{equation}\label{eq:sharpconst}
\sup_{h\in\mathcal V,\ \|h\|_{L^{2}(\pi)}\le\delta}\ \tfrac12\,\big|\big\langle\mathcal Bh,h\big\rangle_{L^{2}(\pi)}\big|
\;=\;\tfrac12\,\delta^{2}\,\big\|P_{\mathcal V}\mathcal B_{\mathrm{sym}}P_{\mathcal V}|_{\mathcal V}\big\| ,
\end{equation}
a supremum that need not be attained.
\item[(ii)] \emph{Bracketing.} The constant in \eqref{eq:sharpconst} obeys
\begin{equation}\label{eq:bracket}
\big\|\vartheta'\sigma_q/\pi\big\|_\infty
\;\le\;\big\|P_{\mathcal V}\mathcal B_{\mathrm{sym}}P_{\mathcal V}|_{\mathcal V}\big\|
\;\le\;\big\|\mathcal B_{\mathrm{sym}}\big\|\;\le\;\big\|\mathcal B\big\| ,
\end{equation}
where the first inequality holds under O-16\textup{(b)} and the outer two hold always. \emph{The
left-hand member is a lower bound and is not the sharp constant; the right-hand members are upper bounds
and are not the sharp constant either.}
\item[(iii)] \emph{Positivity, and it is a hypothesis.} Under
\textbf{O-16\textup{(a)}} --- there exists $h\in\mathcal V$ with $\vartheta'E[q\,(\Delta h)^{2}]\ne0$ ---
the constant in \eqref{eq:sharpconst} is strictly positive, so
$\sup_{\|h\|\le\delta}\tfrac12|\langle\mathcal Bh,h\rangle|\asymp\delta^{2}$ with a constant free of
$\delta$, and no bound of the form $o(\delta^{2})$ as $\delta\downarrow0$ holds uniformly over
$\{h\in\mathcal V:\|h\|_{L^{2}(\pi)}\le\delta\}$. Under \textbf{O-16\textup{(b)}} ---
$\vartheta'\sigma_q\not\equiv0$ --- the constant is at least $\|\vartheta'\sigma_q/\pi\|_\infty>0$, and
O-16\textup{(b)} implies O-16\textup{(a)}.
\item[(iv)] \emph{No weak-norm bound.} Under O-16\textup{(b)} there is no finite $\lambda$ with
\begin{equation}\label{eq:noweak}
\tfrac12\big|\big\langle\mathcal Bh,h\big\rangle_{L^{2}(\pi)}\big|
\;\le\;\lambda\,\big\|\mathcal A_{\mathrm{cond}}h\big\|^{2}_{L^{2}(\Pi)}
\qquad\text{for all }h\in\mathcal V ;
\end{equation}
nor with $\|\mathcal A_{\mathrm{cond}}h\|^{2}$ replaced by
$\|h\|_{L^{2}(\pi)}\,\|\mathcal A^\ast_{\mathrm{cond}}\mathcal A_{\mathrm{cond}}h\|_{L^{2}(\pi)}$.
\end{enumerate}
\end{proposition}

\begin{definition}[The estimator: iterated linearisation with sieve Landweber--Fridman]\label{def:lfstep}
Let $\hat\varphi^{(0)}$ be a preliminary estimator which is a \emph{measurable random element of
$\mathcal V$ with respect to the sample it is computed on} --- written $\sigma(\mathcal I_\ast)$ below, with
$\mathcal I_\ast$ fixed in the verification paragraph, and \emph{this is a requirement on the preliminary and
is stated \textbf{here}, where the preliminary is declared, rather than only inside the verification that
consumes it}; (P-a.1) of \cref{prp:phirate} carries the same requirement, and it is of the kind read off
a formula, which is why it costs nothing to impose --- \textbf{and which lies in $\mathcal S_n$}, the
sieve fixed at the foot of this paragraph; and let $\hat\varepsilon_0\to0$ be a \emph{known}
deterministic sequence with
$P\big(\|\hat\varphi^{(0)}-\vdagger\|_{L^{2}(\pi)}\le\hat\varepsilon_0\big)\to1$ and
$P\big(\|\hat\varphi^{(0)}-\vdagger\|_{H^{s}}\le\mathsf R\big)\to1$ --- the second being a statement about
a finite quantity only because $\hat\varphi^{(0)}\in\mathcal S_n$ and \eqref{eq:sieveinhs} places
$\mathcal S_n$ in $H^{s}$, while $\vdagger\in H^{s}$ by H1's source set.
\emph{The sieve-membership clause is \textbf{imposed here}}, and (P-a.1) of \cref{prp:phirate}, this
opening's companion, carries it in terms. \emph{The clause is not idle}: every iterate after the zeroth
lies in $\mathcal S_n$
% ARCHAEOLOGY.  Through dispatch 20 this opening required measurability into V only
% and then imposed the H^s probability on the strength of a membership stated
% nowhere in it; the fourth bullet below establishes measurability of
% ||hat-varphi^{(j)} - varphi_dagger||_{H^s} expressly "at the configuration
% lem:pathcont states, where hat-varphi^{(0)} in S_n", a CONDITIONAL discharge of
% an UNCONDITIONAL requirement.  No result's hypothesis set changed.
by construction, so without it the preliminary is the one element of the path that need not, and the
$H^{s}$ display above is then a statement about a quantity no clause makes finite; \emph{envelopes, not $O_P$
statements, and \cref{rem:paenvelope} is why}: the constants $\mathsf c_{14},\mathsf c_{15}$ below are
printed before the data are seen, and an $O_P$ hypothesis supplies no number for them to dominate; let
$\mathcal S_n:=\operatorname{span}(b_1,\dots,b_{L_x})\cap\mathcal V$ with $b_1,\dots,b_{L_x}$
the B-spline basis of A8. \emph{Normalise the scale.} Let
\begin{equation}\label{eq:normconst}
\omega_n\;:=\;2\,(2c_C)^{1/2}\,\mathsf c_2^{2}\,c_\Pi\,\|1/\pi\|_\infty^{1/2}\,h_u^{-1}h_x^{-1/2}
\end{equation}
be the bound of \cref{prp:hatAprop}(i), which is \emph{deterministic} and holds \emph{surely} at every
$\varphi\in\mathcal V$, and run the iteration below on $\omega_n^{-1}\hat{\mathcal A}_\varphi$ and
$\omega_n^{-1}\hat T_{h_x}$, so that $\|\omega_n^{-1}\hat{\mathcal A}_\varphi\vert_{\mathcal S_n}\|\le1$
surely. \emph{The letter is $\omega$ because \cref{lem:scaleinv}(i) shows the normalisation \textbf{is} a
Landweber relaxation parameter, which is what $\omega$ names in the regularisation literature.}
% ARCHAEOLOGY.  The letter is not varsigma, which is a Sec.1.6 local bound variable
% and one subscript from the binding varsigma_0 of (P-a.2).  Glossary C67.
By \cref{lem:scaleinv}(i)--(ii) this is a choice of relaxation parameter and changes neither the
iterates nor the stopping index; by (iii)--(iv) it leaves the rate \eqref{eq:phirate} unchanged; by (v)
it inflates the stopping index by $\omega_n^{2}$, and the iteration cap is sized accordingly,
\begin{equation}\label{eq:cap}
\bar m_n\;\ge\;\omega_n^{2}\,\hat\delta_n^{-2} .
\end{equation}
% ARCHAEOLOGY.  Through round 9 the normalisation was asserted to "cost nothing"
% and the cap was delta-hat_n^{-2}; eq:cap is what it actually costs, and it is a
% computation and not an assumption.

\emph{The stopping levels are computed before the data are touched.} Fix $\mathsf c_{14},\mathsf c_{15}<\infty$
and $\kappa:=2\mu/(2\mu+1)$, and define the deterministic \emph{stagewise noise schedule}
\begin{equation}\label{eq:lfsched}
\hat\delta^{\mathrm{eff}}_{n,k}\;:=\;\mathsf c_{14}\Big(\hat\delta_n\ln n\;+\;h_u^{-2}\,\hat\varepsilon_{k-1}^{2}\Big),
\qquad
\hat\varepsilon_{k}\;:=\;\mathsf c_{15}\big(\hat\delta^{\mathrm{eff}}_{n,k}\big)^{\kappa},
\qquad k\ge1 .
\end{equation}
\emph{Every quantity in \eqref{eq:lfsched} is \textbf{computable} by the analyst, from four inputs and
two constants, and the two constants are not free}: $\hat\delta_n$ and $h_u$ are already inputs to the
procedure, $\kappa$ is a function of $\mu$, and $\hat\varepsilon_0$ is the declared rate of the
preliminary --- those are the four; $\mathsf c_{14}$ and $\mathsf c_{15}$ are the two, and they must be
\emph{chosen}, because they are not identified by the data. \emph{What they must satisfy is this.}
% ARCHAEOLOGY.  A round-9 draft said instead that every quantity in eq:lfsched was
% "known".
\Cref{prp:phirate} Step~2
(L3) needs $\mathsf c_{14}\ge\mathsf c_{16}$, with $\mathsf c_{16}$ the constant of \eqref{eq:effnoise};
Step~3 needs $\mathsf c_{15}\ge\mathsf c_{17}$, with $\mathsf c_{17}=2c_\mu\mathsf R_\ast^{1/(2\mu+1)}$
the constant of \eqref{eq:linlf}. Neither $\mathsf c_{16}$ nor $\mathsf c_{17}$ is known to the analyst:
the first is built from the kernel constants and the design, the second from (F-f)'s $c_\mu$ and the
\emph{spectral pre-image} radius $\mathsf R_\ast$ of (P-a.3).
\emph{The radius here is $\mathsf R_\ast$ and \textbf{not} $\mathsf R$, and the two must not be the same
symbol --- this is the same species of defect as $\mathsf c_{14}$ versus $\mathsf c_{16}$ and it sits on
the same inequality.} $\mathsf R$ is an \emph{$H^{s}$ ball radius}: it is what \cref{def:srcball}
measures, what A5(vi$'$) and (P-a.1) bound, and what (P-c), (P-g) and \eqref{eq:entropy} price, so it
carries units of Sobolev norm and enters \emph{approximation} statements. $\mathsf R_\ast$ is a
\emph{spectral pre-image norm} in $L^{2}(\pi)$: it is (P-a.3)'s bound on $\varsigma^{(k)}$, it enters
only through (F-f), and it is what $\mathsf c_{15}$ must dominate. \emph{They are not comparable in
either direction without the interpolation step \cref{rem:pasrc}\textup{(1)} shows is unavailable.}
\emph{The two directions of error are not symmetric.} Chosen too large, the level is
conservative: the rule stops earlier, the rate in \eqref{eq:phirate} is unchanged in order and the
constant degrades. \emph{Chosen too small, the implemented level sits below the true data error
at some stage, (L3) fails, and with it (F-f): the discrepancy guarantee is void and the estimator is
whatever the cap \eqref{eq:cap} returns.} \emph{This is the one place where the procedure carries a tuning constant the
% ARCHAEOLOGY.  Through round 11 one symbol carried both R and R_star, on the two
% sides of an inequality this definition has to establish; glossary C69.  The
% too-small failure mode is the one the fixed level tau_DP delta-hat_n ln n of
% round 9 produced, displaced from the schedule's SHAPE to its SCALE.
theory does not pin, and it is said here rather than left to be discovered}; probe CP-13's
sensitivity arm is what measures how much it matters. For $k=1,\dots,\bar k$:
\begin{enumerate}
\item[(a)] form the linear equation in $v\in\mathcal S_n$,
$\ \hat{\mathcal A}_{\hat\varphi^{(k-1)}}v=-\hat T_{h_x}\big(\hat\varphi^{(k-1)},\theta\big)$;
\item[(b)] run Landweber--Fridman on it from $v_0=0$, at the relaxation \eqref{eq:normconst},
$v_{m+1}=v_m-\omega_n^{-2}\,\hat{\mathcal A}^\ast_{\hat\varphi^{(k-1)}}\big[\hat{\mathcal A}_{\hat\varphi^{(k-1)}}v_m
+\hat T_{h_x}(\hat\varphi^{(k-1)},\theta)\big]$, projected onto $\mathcal S_n$ --- which by
\cref{lem:scaleinv}(i) \emph{is} the unrelaxed iteration for the normalised pair;
\item[(c)] stop at
\[
\hat m^{(k)}\;:=\;\min\Big\{m\le\bar m_n:\ \big\|\hat{\mathcal A}_{\hat\varphi^{(k-1)}}v_m
+\hat T_{h_x}(\hat\varphi^{(k-1)},\theta)\big\|_{L^{2}(\Pi)}\;\le\;
\tau_{\mathrm{DP}}\,\hat\delta^{\mathrm{eff}}_{n,k}\Big\},
\]
the discrepancy principle \emph{at the stage-$k$ level} \eqref{eq:lfsched}, \emph{with the convention
that $\hat m^{(k)}:=\bar m_n$ when the set is empty};
\item[(d)] set $\hat\varphi^{(k)}:=\hat\varphi^{(k-1)}+v_{\hat m^{(k)}}$.
\end{enumerate}
Put $\hat\varphi:=\hat\varphi^{(\bar k)}$.

\emph{The output is a measurable random element, and that is \textbf{verified} here rather than required
elsewhere.} \Cref{def:crossfit}\textup{(c$'$.1)} needs each $\hat\varphi^{(-k)}$ to be a
$\sigma(\mathcal I_{-k})$-measurable random element of $\mathcal V$, and clause \textup{(b)} of that
definition names \emph{this} estimator as the object --- so the property cannot be imposed there as a
construction requirement. \emph{One cannot require a property of an object one has already built}, which is
\cref{rem:obligations}'s test read on an estimator rather than on an assumption. \emph{And the discharge
available for the other nuisances is not available here}: A4(iii) and \cref{def:crossfit}\textup{(a)} settle
$\hat\varrho$ ``by inspection of the formula'', which every kernel and every series construction passes;
inspection of the formula settles nothing for a Landweber--Fridman iterate stopped by a discrepancy rule,
\emph{because the stopping index of \textup{(c)} is random}. The verification is three steps, each
elementary and each consuming a feature of the construction a reader could take for bookkeeping, followed by
a fourth reading that this paragraph's own consumers force.
Write $\mathcal I_\ast$ for the fold the iteration is run on --- $\mathcal I_{-k}$ under
\cref{def:crossfit}\textup{(b)}, $\mathcal I_m$ under \cref{def:hatTexp} --- and read every ``measurable''
below at $\sigma(\mathcal I_\ast)$.
% ARCHAEOLOGY.  The letter was cal J through dispatch 21; the change is dispatch
% 22's.  cal J is the SINGULAR-SYSTEM INDEX SET of prp:compact's statement,
% cal J subset N, so one calligraphic letter carried an observation index here
% (i in cal J) and a singular-value index there (j in cal J), both in STATEMENTS,
% neither inside a proof.  The starred cal I is not a new letter: cal I is already
% this project's letter for a fold (cal I_a, cal I_m, cal I_k, cal I_{-k}), this
% object IS one of those two folds, and the star is the corpus's mark for the
% designated member of a family (cal T_star, cal G_star, Xi_star, sans-R_star,
% omega_star).  Glossary C81.
\emph{One index warning, because both readings occur in this paragraph}: the
superscript in $\hat\varphi^{(-k)}$ and $\mathcal I_{-k}$ just above is \cref{def:crossfit}'s \emph{fold}
index; every $k$ from here to the end of the verification is this definition's \emph{stage} index, and the
two are one glyph apart.
\begin{enumerate}
\item[--] \emph{The iterates.} Suppose $\hat\varphi^{(k-1)}$ is a $\sigma(\mathcal I_\ast)$-measurable random
element of $\mathcal V$. The finitely many numbers $\Delta_{h_x}\hat\varphi^{(k-1)}_i$, $i\in\mathcal I_\ast$, are
then random variables, being values of the \emph{bounded linear} functionals \eqref{eq:smdiff} at a measurable
random element --- \emph{and the step from a random function to a random number is the
\textbf{Carath\'eodory} criterion and not the evaluation of a fixed functional, because the functional
\eqref{eq:smdiff} is itself \textbf{random}}, its kernel carrying $X_{i1},X_{i2}$ and $s_{0,i1},s_{0,i2}$.
``Value of a bounded linear functional at a measurable random element'' gives measurability only for a
\emph{fixed} functional; what is available, and what is needed, is the superposition argument: the map
$F_i(\omega,\varphi):=\Delta_{h_x}\varphi_i(\omega)$ is $\sigma(\mathcal I_\ast)$-measurable in $\omega$ at each
fixed $\varphi$, because $X_{i1},X_{i2}$ are random variables, $s_0$ is Borel by \cref{lem:intker}(iv) and
$C_{h_x}$ is jointly Borel, so the $dx$-integral in \eqref{eq:smdiff} is measurable in $\omega$ by Tonelli
\emph{applied to positive and negative parts} --- \emph{and the qualification is not decoration here}: the
integrand of \eqref{eq:smdiff} is a difference of two boundary-corrected kernels of order $p_K$, which
\eqref{eq:truncmom} makes \emph{sign-changing} by construction, against a $\varphi$ that integrates to zero
against $\pi$, so Tonelli in its nonnegative form does not apply to it as written. The split is legitimate
because the integrand is absolutely integrable at every sample point: $|C_{h_x}|\le\mathsf c_2$ by C3(ii)
and \textup{(P9)}, $\mathcal X$ is a bounded interval by A0, and
$\int_{\mathcal X}|\varphi|\,dx\le|\mathcal X|^{1/2}\|1/\pi\|^{1/2}_\infty\|\varphi\|_{L^{2}(\pi)}<\infty$
by C1(ii) and the Cauchy--Schwarz inequality, so each of the four one-signed integrals is finite and their
signed combination is well defined and measurable. \emph{This is the treatment the same paragraph gives the
$\Pi$-pairing below}; it is also
A2c\textup{(m.2)}'s and \cref{def:crossfit}\textup{(c$'$.1)}'s, one Part away on the other
side of the pairing. The integral is, moreover,
$\|\cdot\|_{L^{2}(\pi)}$-continuous in $\varphi$ at each fixed $\omega$, which is exactly the boundedness of
\eqref{eq:smdiff} that \cref{def:hatTexp} records and which holds at every sample point. Since $\mathcal V$
is a subspace of $L^{2}(\mathcal X,\pi)$ with $\mathcal X$ compact, it is separable, so the Carath\'eodory
criterion makes $F_i$ jointly $\sigma(\mathcal I_\ast)\otimes\mathfrak B(\mathcal V)$-measurable, and $F_i$ then
composes with the measurable map $\omega\mapsto\big(\omega,\hat\varphi^{(k-1)}(\omega)\big)$. \emph{The
device is the corpus's own and is one Part away}: it is what \cref{def:crossfit}\textup{(c$'$.1)} names for
\eqref{eq:prodmeas}, what the opening of the proof of \cref{prp:frechet} runs for $G_\varphi$ and its
$u$-derivatives, and what C2$'$(i) runs for the envelopes $\bar f_k$;
% ARCHAEOLOGY.  Through dispatch 19 this link was stated as an evaluation, which is
% the one reading the randomness of eq:smdiff forecloses.
and $\hat T_{h_x}(\hat\varphi^{(k-1)},\theta)$, $\hat{\mathcal A}_{\hat\varphi^{(k-1)}}$ and
$\hat{\mathcal A}^\ast_{\hat\varphi^{(k-1)}}$ depend on $\hat\varphi^{(k-1)}$ through those numbers
\emph{alone}, by \eqref{eq:hatTx}, \eqref{eq:hatg} and \eqref{eq:hatAcon}, the adjoint being determined by the
operator and the two inner products of \cref{def:spaces}. Restricted to $\mathcal S_n$, which
is a \emph{fixed} finite-dimensional space, the iteration \textup{(b)} is therefore an affine recursion on
a Euclidean space of that dimension whose coefficients are random variables, started at the constant
$v_0=0$; so each
$v_m$, $m=0,1,\dots,\bar m_n$, is a $\sigma(\mathcal I_\ast)$-measurable random element of $\mathcal S_n$, being a
finite composition of measurable maps. \emph{That the coefficients are random variables is
\textbf{Tonelli} and not composition, and the distinction is the same one the previous link turns on.}
\emph{The coordinates are those of a basis of $\mathcal S_n$ and \textbf{not} of $\{b_1,\dots,b_{L_x}\}$},
and two facts force the distinction:
$\mathcal S_n=\operatorname{span}(b_1,\dots,b_{L_x})\cap\mathcal V$, and the linear functional
$\zeta\mapsto\int_{\mathcal X}\zeta\,d\pi$ is \emph{not} identically zero on
$\operatorname{span}(b_1,\dots,b_{L_x})$: A8's $b_j$ are B-splines, hence \emph{nonnegative} and not
identically zero, and $\pi$ is bounded below away from zero by C1(ii) --- so
$\int b_j\,d\pi>0$ for every $j$, no $b_j$ lies in $\mathcal V$, $\{b_j\}$ spans a strictly larger space, and
$\dim\mathcal S_n=L_x-1$. \emph{Both consequences matter}: the coordinate map written in the $b_j$ is not a
coordinate map on $\mathcal S_n$, and $\hat{\mathcal A}_{\hat\varphi^{(k-1)}}b_j$ evaluates the operator off
the domain \cref{def:spaces} declares for it. \emph{Neither touches the conclusion}, and that is why the
repair is a rewording: fix once any basis of $\mathcal S_n$ --- every element of it lies in $\mathcal V$, so
the operator is applied only where it is defined, and every element is a finite linear combination of
$b_1,\dots,b_{L_x}$ --- and read the recursion in its coordinates. Its matrix is that of
$\hat{\mathcal A}^\ast_{\hat\varphi^{(k-1)}}\hat{\mathcal A}_{\hat\varphi^{(k-1)}}\vert_{\mathcal S_n}$,
whose entries are the $L^{2}(\Pi)$ pairings of the images under
$\hat{\mathcal A}_{\hat\varphi^{(k-1)}}$ of two basis elements, and its constant vector carries the
$L^{2}(\Pi)$ pairings of those images with $\hat T_{h_x}(\hat\varphi^{(k-1)},\theta)$
--- each an \emph{integral over $(u,t)$ against $\Pi$ of a random integrand}, \eqref{eq:hatAadj} writing the
adjoint as exactly such an integral. \emph{A finite composition of measurable maps does not reach an
integral}; what does is joint measurability in $(\omega,u,t)$ followed by Tonelli's theorem. The joint
measurability is again Carath\'eodory: by \eqref{eq:hatAcon} and \eqref{eq:hatg}, at a fixed element $\zeta$
of $\mathcal V$ the maps
$(u,t)\mapsto\hat{\mathcal A}_{\hat\varphi^{(k-1)}}\zeta(u,t)$ and
$(u,t)\mapsto\hat T_{h_x}(\hat\varphi^{(k-1)},\theta)(u,t)$ are continuous at every sample point --- $C_{h_u}$
and $\mathsf C_{h_u}$ being continuous in $u$ by \cref{lem:intker} and $t\mapsto w(t'\Phi(\Delta Z_i))$
continuous --- while each is $\sigma(\mathcal I_\ast)$-measurable in $\omega$ at fixed $(u,t)$ by the paragraph
above, so both are $\sigma(\mathcal I_\ast)\otimes\mathfrak B(\IR\times\mathcal T)$-measurable; each factor lies
in $L^{2}(\Pi)$ --- by \cref{prp:hatAprop}(i) for the operator and by $\mathsf c_3c_\Pi<\infty$ for
$\hat T_{h_x}$ --- so the products are $\Pi$-integrable by the Cauchy--Schwarz inequality, and Tonelli
applied to positive and negative parts then makes each pairing a $\sigma(\mathcal I_\ast)$-measurable function
of $\omega$. \emph{This is the inheritance
\cref{def:crossfit}\textup{(c$'$.1)} performs for $\omega\mapsto\|\hat\ell^{(-k)}-\ell_\star\|_{L^{2}(P)}$,
read on $\Pi$ instead of $P$.}
\item[--] \emph{The stopping index.} $\tau_{\mathrm{DP}}$ and $\hat\delta^{\mathrm{eff}}_{n,k}$ are
deterministic by \eqref{eq:lfsched}, and $\|\hat{\mathcal A}_{\hat\varphi^{(k-1)}}v_m+\hat T_{h_x}(\hat\varphi^{(k-1)},\theta)\|_{L^{2}(\Pi)}$
is a random variable for each $m$ --- \emph{by the previous step \textbf{and Tonelli}, the two being
needed and neither sufficing alone}, since an $L^{2}(\Pi)$ norm is a $\Pi$-integral of a random function
and not a value of one: the previous step makes
$(\omega,u,t)\mapsto\big(\hat{\mathcal A}_{\hat\varphi^{(k-1)}}v_m+\hat T_{h_x}(\hat\varphi^{(k-1)},\theta)\big)(u,t)$
jointly measurable, $v_m$ being $\sigma(\mathcal I_\ast)$-measurable with values in the fixed
finite-dimensional $\mathcal S_n$, and Tonelli then makes $\omega\mapsto\int(\cdot)^{2}\,d\Pi$ measurable, the square root
being continuous. So, for each $m\le\bar m_n$, the event that \emph{that}
norm is at most $\tau_{\mathrm{DP}}\hat\delta^{\mathrm{eff}}_{n,k}$ is measurable --- \emph{it carries no
letter}, every short candidate being claimed --- and $\{\hat m^{(k)}=m\}$ is that event intersected with the
complements of the same events at the indices $m'<m$: a \emph{finite} intersection of measurable sets.
\emph{Two further features of \textup{(c)} --- not the two the paragraph after this one names, which are
about the \emph{level} --- are what this step consumes}: the cap
$\bar m_n$ of \eqref{eq:cap} is \emph{deterministic}, so the family is \textbf{finite} and no limit is taken;
and the emptiness convention makes $\hat m^{(k)}$ \textbf{total}, so no sample point is left without a value
and $\{\hat m^{(k)}=\bar m_n\}$ absorbs the residual event. \emph{A random cap or a partial index would break
the argument at this line and nowhere else.}
\item[--] \emph{The output.} $v_{\hat m^{(k)}}=\sum_{m\le\bar m_n}\un\{\hat m^{(k)}=m\}\,v_m$ is a
\emph{finite} sum of indicators of measurable events against measurable random elements, hence a
$\sigma(\mathcal I_\ast)$-measurable random element of $\mathcal S_n$; and $\hat\varphi^{(k)}$ of \textup{(d)} is
one of $\mathcal V$, $\mathcal V$ being a linear space and $\mathcal S_n\subset\mathcal V$
(\cref{def:spaces}). Induction on $k\le\bar k$, from a $\sigma(\mathcal I_\ast)$-measurable
$\hat\varphi^{(0)}$ --- \emph{which is a requirement on the preliminary, and \textbf{is} of the kind read off a
formula; it is \textbf{stated} in this definition's opening and in (P-a.1) of \cref{prp:phirate}, and is
cited rather than imposed here} --- gives the claim for $\hat\varphi=\hat\varphi^{(\bar k)}$.
\item[--] \emph{The $H^{s}$ norm of the same object, which the paragraph owes because \cref{lem:pathcont} and
A5(vi$'$) are $H^{s}$ statements.} Both take probabilities of
$\max_{j\le\bar k}\|\hat\varphi^{(j)}-\vdagger\|_{H^{s}}\le\mathsf R$, so
$\omega\mapsto\|\hat\varphi^{(j)}-\vdagger\|_{H^{s}}$ must be a random variable as well, and measurability
into the $L^{2}(\pi)$ container does not by itself give it. \emph{It follows from the three steps above at the
configuration \cref{lem:pathcont} states}, where $\hat\varphi^{(0)}\in\mathcal S_n$ and hence the whole path
lies in $\mathcal S_n$: since $\mathcal S_n\subset H^{s}$ by \eqref{eq:sieveinhs}, $\|\cdot\|_{H^{s}}$ is a
norm on a \emph{fixed} finite-dimensional space, hence a continuous --- and so measurable --- function of the
coordinates, which the steps above make random variables. \emph{The membership is A8's business and it is now
\textbf{stated} there rather than flagged here}: nothing in this definition places the B-spline basis in
$H^{s}$, and \eqref{eq:sieveinhs} is A8's basis clause read as the containment it delivers.
\emph{\eqref{eq:jackson} is \textbf{not} the right pointer, and the error would be one of direction
rather than of precision}: \eqref{eq:jackson} states how well splines \emph{approximate}
$H^{s}$ functions, and what is needed is that splines \emph{lie in} $H^{s}$. \emph{Approximation power does
% ARCHAEOLOGY.  Through dispatch 19 this bullet pointed at eq:jackson.
not give membership, and the two conditions are genuinely different inequalities on the same degree} --- a
basis of degree $d$ approximates $H^{s}$ at the rate \eqref{eq:jackson} states as soon as $s\le d+1$, while
$\mathcal S_n\subset H^{s}$ requires $d>s-\tfrac12$; on the window $s-1\le d\le s-\tfrac12$ the first holds
and the second fails, and the witness $s=1.7$, $d=1$ sits in it. A8's basis clause is the
right home, as this bullet's own first sentence already said, and it is stated there in \emph{degree}, the
one reading of ``order'' on which the spline conventions agree.
\end{enumerate}
\emph{This is a discharged verification and \textbf{not} a fifth entry in \cref{rem:obligations}'s inventory,
and the distinction is that remark's own.} That inventory carries obligations whose \emph{discharge site is a
later step} --- something a stated result's truth depends on and no present result establishes. Here the
discharge is performed in the same definition that constructs the object, so nothing is owed forward and there
is no consumer to name; the precedents are A4(iii)'s joint-measurability half and
\cref{def:crossfit}\textup{(c$'$.1)}'s Tonelli inheritance, neither of which is inventoried either, for the
same reason. \emph{Three stated results consume this paragraph}, each a probability statement about
$\hat\varphi$ and none of them well posed without it: \eqref{eq:phirate}; \cref{lem:pathcont}'s bound on
$\max_{j\le\bar k}$ with probability tending to one; and (P-a.3)'s event, which is defined through the realised
operator at $\hat\varphi^{(k-1)}$.

\emph{Two features of (c) are load-bearing and one of them is a correction.} \emph{First, the level is
stage-dependent, and it must be.} The data error of the stage-$k$ equation (a) is not $\hat\delta_n$
but $\hat\delta_n\ln n$ \emph{plus} the linearisation remainder $h_u^{-2}\varepsilon_{k-1}^{2}$, where
$\varepsilon_{k-1}=\|\vdagger-\hat\varphi^{(k-1)}\|_{L^{2}(\pi)}$ --- see \eqref{eq:stagedecomp}. A
fixed level $\tau_{\mathrm{DP}}\hat\delta_n\ln n$ sits
\emph{below} the true data error in exactly the regime the recursion must pass through, namely
$h_u^{-2}\varepsilon_{k-1}^{2}\ge\hat\delta_n\ln n$; there the classical discrepancy guarantee does not
apply and the rule either over-fits the noise or fails to stop before $\bar m_n$.
\emph{Second, the level is deterministic and supplied from outside, and the reason has two parts of
very different status. The first is this paper's own and is sharp; the second is classical and cited,
and is weaker than it looks.}
% ARCHAEOLOGY.  Rounds through 9 implemented the fixed level tau_DP delta-hat_n ln n,
% a mismatch between the estimator and its analysis; it is repaired at the estimator
% rather than in the proof.  A round-9 draft over-read Bakushinskii's veto twice.

\emph{(2a) The realised residual cannot supply the level, and the reason is what ill-posedness is.} The
natural data-driven repair is to replace $h_u^{-2}\hat\varepsilon_{k-1}^{2}$ by
$h_u^{-2}\hat r_{k-1}^{2}$ with $\hat r_{k-1}$ the residual realised at stage $k-1$. \emph{It cannot
work.} The quantity the level must dominate is $h_u^{-2}\varepsilon_{k-1}^{2}$ with
$\varepsilon_{k-1}=\|\vdagger-\hat\varphi^{(k-1)}\|_{L^{2}(\pi)}$, and a residual bounds $\varepsilon$
from below only through a lower bound on $\|\hat{\mathcal A}_{\varphi_\circ}v\|$ in terms of $\|v\|$.
\emph{That $\|\hat{\mathcal A}v_\dagger\|\ll\|v_\dagger\|$ is available is not an accident of this
construction; it is the definition of ill-posedness.} So $\hat r_{k-1}$ can be arbitrarily smaller than
$\varepsilon_{k-1}$, and a threshold built from it sits below the data error for exactly the reason the
fixed threshold did. \emph{This observation is the paper's own and it is the sharp part of the
argument.}

\emph{(2b) The general statement is Bakushinskii's veto, it is cited, and it is weaker than it looks in
a statistical setting.} That \emph{no} rule depending on the data alone --- and not on the noise level
--- can make a regularisation method convergent for a genuinely ill-posed problem is
\citet{bakushinskii1984}; see \citet[\S3.3]{engl2000}. \emph{Two scope statements, because the veto is
routinely over-read.} It is a \emph{worst-case, deterministic}
statement: convergence is required uniformly over the admissible data, and the counterexample it turns
on is a worst case. \emph{In the statistical setting this paper inhabits it does not rule out
data-driven selection}: Lepskii-type balancing over the regularisation index is purely data-driven and
attains the optimal rate up to a logarithmic factor, and it is a device this project already runs
elsewhere. \emph{So the schedule is preferred here, not forced here}, and the reasons are three and
local. The number of stages $\bar k$ is finite and $\bar k_a$ is deterministic, so there is no
adaptation problem across stages to solve. A deterministic level is what makes the rule
\cref{def:lfstep}(c) implements and the rule \cref{prp:phirate} Step~2 analyses coincide \emph{by
construction}. And a Lepskii stop over $m$ would
need a variance proxy at each candidate index that the operator-noise model of this paper does not
supply without a further bound. \emph{A balancing stop is therefore a live alternative rather than a
foreclosed one, it is registered as the third arm of probe CP-13, and if it matches, the requirement
that (P-a.2) deliver a \emph{rate} --- open item Q-1 --- softens materially.}

\emph{The price actually paid, which is separate from both.} The schedule requires a \emph{rate}
$\hat\varepsilon_0$ for the preliminary and not mere consistency. \emph{That price is not new, only
newly visible}: the claim that $\bar k$ is finite and does not grow with $n$ already required
$\hat\varepsilon_0$ to be polynomially small. Hypothesis (P-a.2) of \cref{prp:phirate} carries it.

\emph{Every Landweber iteration in this construction is applied to a \textbf{linear} equation}; the
nonlinearity is carried by the outer index $k$ and by nothing else. That is the entire difference from
the route \cref{rem:lfnonlinear} declines, and it is what makes a cone condition unnecessary.
\end{definition}

\begin{proposition}[\texttt{prp:phirate}: the rate for the $\varphi$-step]\label{prp:phirate}
Suppose A0, A1, A1b(iii)--(iv), A4, A5(vi$'$), B1(i)--(v), C1(ii), C2$'$, C3 and C5$'$(i) hold.
\emph{Three of those are narrowed from the blocks ``B1'', ``C1'' and ``C5$'$'' that stood here, and
each narrowing removes a clause this proposition cannot have}: B1(vi) is \textbf{retired} and is not available
as a hypothesis of any result, so a block citation of B1 charged this proposition with it; C1(i) is
consumed at exactly one site in the corpus, Step~2 of \cref{lem:aerr}, the route this proposition's Step~0
\textbf{declines}, while C1(ii) is consumed here transitively through \cref{prp:hatAprop} and
\cref{lem:nusigma}; and C5$'$(ii) is a \emph{lower} bound on $h_x$ while this proposition's own (P-g) is an
\emph{upper} bound on $h_x$, so a block citation of C5$'$ put a floor and a ceiling in one hypothesis set ---
\emph{and, at the inherited strengthening of C5$'$(ii), a floor above the ceiling}. \emph{The two do not meet
at every admissible parameter, and the failure is not a corner}: at $(\mu,s,a,r,d_Z)=(2,1.2,0.3,\tfrac83,3)$,
which clears \eqref{eq:frontier} at $\kappa\alpha=0.256$ and clears \eqref{eq:lxwindow} at $s>1$, the
strengthened floor is $h_x\gg n^{-0.44}$ and (P-g) is $h_x=O(n^{-0.4571})$. The threshold is
\eqref{eq:hxwindow} and the emptiness band is non-empty at every $\alpha\in(\tfrac14,\approx0.348)$, which
\eqref{eq:frontier} does not exclude. \emph{\textbf{What follows from that is vacuity, not falsity}}: on the
block-cited reading this proposition's hypotheses were unsatisfiable there and it held on an empty
hypothesis set, and no conclusion of this corpus is false at that configuration. The narrowing removes the
conflict from the hypothesis set; C5$'$(ii)'s own restatement removes it from the maintained block.
Suppose H1 holds, for
the two things this proposition uses it for and \textbf{not} for a third --- \emph{it fixes the Sobolev
scale $H^{s}$ that \cref{def:srcball}, A5(vi$'$), (P-c) and (P-g) are written in, and it supplies A2's
Picard form}, but \emph{the identity $\mu=s/(2a)$ is \textbf{not} assumed and is \textbf{not} used};
$\mu$ is primitive, at (P-a.3), and \cref{rem:pasrc}\textup{(4)} says exactly what pinning it to
$s/(2a)$ would cost. Suppose A8 holds \emph{in the two clauses this proposition consumes and not in the
third} --- its \emph{basis} clause, that $b_1,\dots,b_{L_x}$ is a B-spline basis of degree exceeding
$s-\tfrac12$ on a quasi-uniform mesh over $\mathcal X$, which is what delivers \eqref{eq:sieveinhs} and
with it the well-posedness of (P-a.1) below, and its \emph{upper end} $L_x\ll n^{1/2}$, but \textbf{not}
its lower end $L_x\gg n^{1/(2s)}$, which is the balancing route's and which A8's own comment already
records as unused on the rate route; suppose in addition $p_K\ge s-\tfrac12$, which \cref{lem:unifop}
Step~4 consumes and which, $p_K=8$ being \emph{fixed} by C3 and by C5$'$(i), is a hard cap
$s\le\tfrac{17}2$ and not a free condition. \emph{It is \textbf{not} free at every $s$ the budget table
admits}: $p_K$ is not a parameter the analyst may raise against $s$, so the clause caps $s$, while
% ARCHAEOLOGY.  Through dispatch 20 this line read "which at p_K = 8 is free at
% every s the budget table admits", which is FALSE.
\eqref{eq:lxwindow} below is a \emph{floor} on $s$ that diverges as $\alpha\uparrow\tfrac12$. \emph{The two
are incompatible at every $\alpha\ge\tfrac{68}{145}\approx0.4690$, at every $s$ and every $a$}: the floor
$9\alpha/\big(8(1-2\alpha)\big)$ reaches $\tfrac{17}2$ there. \emph{Witness, and it is admissible}: $r=8$,
$d_Z=1$, so $\alpha=\tfrac8{17}$; with $\mu=2$ the frontier \eqref{eq:frontier} clears at
$\kappa\alpha=\tfrac{8}{17}\cdot\tfrac45\approx0.3765>\tfrac14$; \eqref{eq:lxwindow} demands $s>9$ and the
cap allows $s\le\tfrac{17}2$; \emph{no $s$ exists}. \emph{The disposition follows A4(ii)'s at
$d_Z=0$ and is the same one}: on $[\tfrac{68}{145},\tfrac12)$ this proposition's hypotheses are
unsatisfiable, so it is \textbf{vacuous rather than false} there, and \cref{thm:rootn}'s group (T-5) is
unsatisfiable with it. A4(ii)'s $d_Z\ge1$ repair delivers only $\alpha<\tfrac12$ and does not
reach this range, which is entirely admissible; the cap is recorded as a row of the smoothness-budget table
and no further restriction on $\alpha$ is imposed here. And suppose
\begin{enumerate}
\item[(P-a)] \emph{(preliminary, and the \textbf{spectral} source condition at the compression)} Three
clauses, all three \emph{envelopes holding with probability tending to one} rather than $O_P$
statements, and the distinction is load-bearing --- see \cref{rem:paenvelope}.
\begin{enumerate}
\item[(P-a.1)] \emph{(preliminary, $H^{s}$ --- and $\sigma(\mathcal I_\ast)$-measurable)}
$\hat\varphi^{(0)}$ is a $\sigma(\mathcal I_\ast)$-measurable random element of $\mathcal V$, in the sense and
with the fold $\mathcal I_\ast$ of \cref{def:lfstep}; $\hat\varphi^{(0)}\in\mathcal S_n$; and
$P\big(\|\vdagger-\hat\varphi^{(0)}\|_{H^{s}}\le\mathsf R\big)\to1$, with $\mathsf R$ the
\emph{same} deterministic radius as in \cref{def:srcball}, A5(vi$'$), (P-c) and (P-g).
\emph{The measurability half is a requirement on the preliminary and is stated \textbf{here} because
some hypothesis list must carry it}: \cref{def:lfstep}'s third bullet inducts on $k$ \emph{from} it, so
the verification that the estimator is a random element consumes it.
% ARCHAEOLOGY.  Through dispatch 19 the measurability half was stated only inside
% that verification -- a hypothesis a proof reads that no hypothesis list stated.
\emph{The two remaining halves are jointly satisfiable only under \eqref{eq:sieveinhs}, and that is not
a technicality}: $\hat\varphi^{(0)}\in\mathcal S_n$ forces $\|\vdagger-\hat\varphi^{(0)}\|_{H^{s}}$ to be
read at a sieve element, so unless $\mathcal S_n\subset H^{s}$ the norm is $+\infty$ surely and the two
halves of this one clause are \emph{jointly unsatisfiable} --- the probability statement would then hold
on the empty event. A8's basis clause supplies \eqref{eq:sieveinhs}, and H1 supplies $\vdagger\in H^{s}$.
\item[(P-a.2)] \emph{(preliminary, rate)} $P\big(\|\vdagger-\hat\varphi^{(0)}\|_{L^{2}(\pi)}\le\hat\varepsilon_0\big)\to1$
for a known deterministic $\hat\varepsilon_0=O(n^{-\varsigma_0})$ whose exponent clears the repelling
fixed point of the stagewise recursion,
\[
\varsigma_0\;>\;\frac{\kappa\,\alpha}{4\,(2\kappa-1)},
\qquad \kappa=\frac{2\mu}{2\mu+1},\quad\alpha=\frac{r}{2r+d_Z}
\]
--- a requirement strictly weaker than the conclusion \eqref{eq:phirate} precisely when
$\mu>\tfrac56$, which is (P-f) \emph{and is the same inequality rather than a second one that agrees
with it}, as Step 4(a.1) of the proof shows.
\item[(P-a.3)] \emph{(source condition, stated \textbf{spectrally} and \textbf{at the compression})}
There are $\mu>0$ and a finite $\mathsf R_\ast$, neither depending on $n$, such that with
\begin{equation}\label{eq:compress}
\hat{\mathcal P}^{\,\mathcal S}_{\varphi}\;:=\;P_{\mathcal S_n}\,\hat{\mathcal A}^{\ast}_{\varphi}\hat{\mathcal A}_{\varphi}\big\vert_{\mathcal S_n}
\;:\;\mathcal S_n\to\mathcal S_n
\end{equation}
the \emph{compression} of $\hat{\mathcal A}^{\ast}_{\varphi}\hat{\mathcal A}_{\varphi}$ to
$\mathcal S_n$,
\[
P\Big(\ \bar v^{(k)}_\dagger=\big(\hat{\mathcal P}^{\,\mathcal S}_{\hat\varphi^{(k-1)}}\big)^{\mu}\varsigma^{(k)}
\ \text{ with }\ \big\|\varsigma^{(k)}\big\|_{L^{2}(\pi)}\le\mathsf R_\ast
\ \text{ for every }k\le\bar k\ \Big)\;\longrightarrow\;1 ,
\]
where $\bar v^{(k)}_\dagger:=P_{\mathcal S_n}\big(\vdagger-\hat\varphi^{(k-1)}\big)$.
\end{enumerate}
\emph{(P-a.3) is \textbf{not} the pair --- a Hilbert-scale source
condition on $\bar v^{(k)}_\dagger$ together with the link condition of \cref{lem:srctransfer}(iii) on
$\mathcal S_n$ --- and does not follow from it}: the pair delivers (P-a.3) only through operator
interpolation at exponent $2\mu$,
which \cref{lem:srctransfer}(ii.a) refutes at \emph{every} $2\mu>1$ while \eqref{eq:muhalf} forces
$2\mu>\tfrac53$, so the pair never implies what (F-f) consumes. \Cref{rem:pasrc} carries the full
disposition. \emph{And the operator is named as the
compression, which is what \eqref{eq:compress} is for}: the operator (F-f) is applied to in Step~2 of
the proof is $\mathcal L^{\ast}\mathcal L$ with
$\mathcal L=\omega_n^{-1}\hat{\mathcal A}_{\hat\varphi^{(k-1)}}\vert_{\mathcal S_n}$, and that
\emph{is} $\omega_n^{-2}\hat{\mathcal P}^{\,\mathcal S}_{\hat\varphi^{(k-1)}}$ and not the restriction
of $\hat{\mathcal A}^{\ast}\hat{\mathcal A}$ to $\mathcal S_n$ --- two operators whose ranges are nested
in neither direction;
% ARCHAEOLOGY.  Rounds 1--11 assumed the Hilbert-scale pair; the replacement is a
% repair and not a strengthening, and is the most consequential amendment of that
% round.  Through round 11 the hypothesis named the restriction of A*A to S_n and
% the proof used the compression.
\item[(P-b)] \emph{\textup{[}struck\textup{]}} --- an operator-estimation budget
$\gamma_n=O(\hat\delta_n^{1/(2\mu+1)})$ belonging
to the perturbation-of-$\mathcal A$ route, which Step~0 of the proof does not take; the label is kept
empty rather than renumbered so that an existing citation of (P-b) lands here;
\item[(P-c)] \emph{(sieve approximation, and it pays for the projection step)}
$\mathsf R\,L_x^{-s}=O\big(h_u\,\hat\delta_n\big)$. \emph{The factor $h_u$ is not decoration}: the
source and link conditions are imposed on $\mathcal S_n$, so the proof works with the sieve projection
$P_{\mathcal S_n}(\vdagger-\hat\varphi^{(k-1)})$ and must carry the residual through
$\hat{\mathcal A}$, whose norm on the complement is bounded only by $O(h_u^{-1})$. \emph{(P-c) is a
\textbf{lower} end for $L_x$, and it is this proposition's own; it \textbf{replaces} A8's lower end
rather than restating it}, being read at $h_u\hat\delta_n$ and not at $n^{-1/2}$;
\item[(P-d)] \emph{(uniformity)} \eqref{eq:entropysieve}, an \emph{upper} end for $L_x$, so that
\cref{lem:unifop} clause (b) is available. \emph{(P-c) and (P-d) are the two ends of one window and
they must be read together}: with $\hat\delta_n=n^{-\alpha}$ and $h_u\asymp\hat\delta_n^{1/8}$ they are
$L_x\gg n^{9\alpha/(8s)}$ and $L_x=o\big(n^{1-2\alpha}/\ln n\big)$, so the window is non-empty if and
only if
\begin{equation}\label{eq:lxwindow}
s\;>\;\frac{9\alpha}{8\,(1-2\alpha)} ,
\end{equation}
which is $s>\tfrac{27}{16}=1.6875$ at $\alpha=\tfrac38$ and $s>1$ at $\alpha=0.32$. \emph{This is a
genuine joint line in the smoothness budget and it is now recorded there.} \emph{Had A8's lower end
been carried as well}, the window would instead require $s>\big(2(1-2\alpha)\big)^{-1}$, which is
$s>2$ at $\alpha=\tfrac38$; that is strictly the stronger line at every $\alpha<\tfrac49$, and it is
why A8 is cited above by clause. A configuration on which \emph{both} routes are to be available at
once must clear the stronger one;
\item[(P-e)] \emph{(stopping)} A7 in the amended form \eqref{eq:a7amend}, with
$\tau_{\mathrm{DP}}>2$;
\item[(P-f)] \emph{(the linearisation remainder is absorbed)}
$\delta_n^{2}=O\big(h_u^{2}\,\hat\delta_n\ln n\big)$, the one bandwidth-explicit condition of Part H;
\cref{rem:h6bandwidth} unpacks it, shows that at the unamended bandwidth $h_u\asymp n^{-1/16}$ it reads
$\tfrac{2\mu-1}{2\mu+1}\cdot\tfrac{r}{2r+d_Z}>\tfrac18$, and shows that at C5$'$(i) as amended ---
$h_u\asymp\hat\delta_n^{1/p_K}$, the largest bandwidth whose bias is $O(\hat\delta_n)$ --- it reads
$\mu>\tfrac56$, \emph{strictly, with the strictness forced by the logarithms};
\item[(P-g)] \emph{($x$-smoothing bias)} $\mathsf R\,h_x^{s-1/2}=O(\hat\delta_n)$, an \emph{upper}
bound on $h_x$, which C5$'$ does not carry;
\end{enumerate}
and suppose finally
\begin{equation}\label{eq:muhalf}
\mu\;>\;\tfrac56 .
\end{equation}
Then there is a finite $\bar k$, not depending on $n$ and computable from $(\mu,\varsigma_0,\alpha)$,
such that the estimator of \cref{def:lfstep} satisfies
\begin{equation}\label{eq:phirate}
\big\|\hat\varphi-\vdagger\big\|_{L^{2}(\pi)}\;=\;O_P(\delta_n),
\qquad
\delta_n\;:=\;\big(\hat\delta_n\ln n\big)^{2\mu/(2\mu+1)} ,
\end{equation}
and $\delta_n=o_P(n^{-1/4})$ whenever \eqref{eq:frontier} holds. \emph{Hence A5(i)--(iii) are
discharged, not assumed. A5(vi$'$) is \textbf{consumed} --- it is a hypothesis of this proposition, at
\cref{lem:pathcont}(c) --- and A5(vi) follows from it at $k=\bar k$}; A5(iv) and A5(v) concern other
objects and are untouched.

\emph{Where the threshold comes from, and it is two conditions and not one.} The stagewise recursion of
Step~4 of the proof is \eqref{eq:fixedpoint},
$\hat\varepsilon_k=\mathsf c_{15}\big(\mathsf c_{14}(\hat\delta_n\ln n+h_u^{-2}\hat\varepsilon_{k-1}^{2})\big)^{2\mu/(2\mu+1)}$,
a recursion in the \emph{deterministic schedule} of \eqref{eq:lfsched} and not in the unobservable
error.
% ARCHAEOLOGY.  The symbol delta_n^{(k)} that a round-9 draft used for the
% unobservable error is withdrawn; glossary C58.
\emph{Contraction} needs $2\kappa>1$ with $\kappa=2\mu/(2\mu+1)$, i.e.\ $\mu>\tfrac12$; \emph{the fixed
point being \eqref{eq:phirate}} needs (P-f), which at the amended bandwidth is $\mu>\tfrac56$. The
second implies the first, so \eqref{eq:muhalf} is stated at $\tfrac56$ and the $\mu>\tfrac12$ line of
the budget table is now slack. \emph{The coincidence worth naming is at $\tfrac12$, not at $\tfrac56$}:
$\mu>\tfrac12$ is also exactly what A5(iii) proves necessary for \eqref{eq:frontier} to be satisfiable
at any $r$, and it is exactly the index above which the inherited nonlinear route's rate theorem is
unavailable (\cref{rem:lfnonlinear}(3)). The two arise independently --- one from the propensity
frontier, one from the linearisation recursion --- and they coincide. \emph{What
\eqref{eq:muhalf} charges beyond that is $(\tfrac12,\tfrac56]$, and that is a real line in the
budget}, listed with the others in the assumption block's table.
\emph{And what happens at $\tfrac56$ is \textbf{not} a second coincidence.} Step 4(a.1) of the proof
exhibits a requirement on the preliminary,
\eqref{eq:prelimrate}, and shows it is weaker than the rate delivered exactly when $\mu>\tfrac56$. That
is (P-f) again --- \emph{the same inequality}, written in the same primitives, up to the strictly
positive factor $\kappa/(2\kappa-1)$. \emph{The
statement worth printing, and it is not a restatement of (P-f) either}: (P-f) is exactly the
% ARCHAEOLOGY.  A round-9 draft called the agreement at 5/6 a coincidence and said
% the two conditions were not restatements of one another; that sentence is struck.
condition under which the newly exposed preliminary-rate requirement is \emph{not self-defeating}, so
that \eqref{eq:prelimrate} adds a genuine new demand on $\hat\varphi^{(0)}$ --- open item Q-1 --- without
making the procedure ask its input to be as good as its output.

\emph{And the load-bearing hypothesis is (P-a.3), which should be read in the same breath as
\eqref{eq:muhalf}.} (P-a.3) is a \emph{random spectral source condition} on a \emph{finite-rank,
sieve-compressed} operator at a \emph{data-dependent} base point, assumed to hold with probability
tending to one, with no sufficient condition supplied anywhere in this file; it is what makes the
conditional-on-sample route legitimate, and it is what carries the source index above $\tfrac12$.
\emph{Any statement that ``the linearised route reaches the same place under the same threshold'' is
true only with (P-a.3) priced in the same sentence} --- see \cref{rem:lfnonlinear},
\cref{rem:pasrc} and open item Q-4.
\end{proposition}

% ============================================================================
%  END OF PART H.
% ============================================================================

% ############################################################################
% ##  PART I -- memo Sec.10 STEP 7.  thm:rootn, at plan Sec.2.6's scope, with
% ##  cor:ratefloor (T2') as the floor.  Dispatch 10.
% ##
% ##  WHAT THIS PART SUPPLIES THAT NO EARLIER PART DID.
% ##   (1) THE MOMENT FUNCTION HAS A DEFINING DISPLAY.  psi is cited in Parts F
% ##       and G (lem:secondord, prp:orthog, prp:sharpquad) and was declared only
% ##       in notation_glossary.md.  def:moment fixes it in the .tex, together
% ##       with J and Sigma, which the memo's target statement names and which no
% ##       file displayed.
% ##   (2) THE ESTIMATOR.  def:crossfit and def:onestep.  def:lfstep defines the
% ##       varphi-step; nothing defined the theta-step.
% ##   (3) THE PRELIMINARY FOR theta.  lem:prelimtheta.  This is the object the
% ##       round-3 strategist-critic's m-23 named, and eq:prelimrate is NOT it --
% ##       see rem:rootnhyp(4).
% ##
% ##  ORDER OF WORK: def:moment, def:crossfit, def:onestep; lem:momzero;
% ##  lem:prelimtheta (L-14 with a rate); lem:biasexp (L-6 and L-7 instantiated
% ##  at the estimated nuisances); lem:epterm (L-8); lem:sandwich (L-11);
% ##  thm:rootn; cor:ratefloor (T2'); prp:medianpart (L-10); three remarks.
% ############################################################################

% ===========================================================================
%  I.1  THE MOMENT FUNCTION, THE JACOBIAN, THE VARIANCE
% ===========================================================================

\begin{definition}[The debiased moment function, and the two matrices]\label{def:moment}
For a nuisance triple $\zeta=(\varphi,L,\varrho)$ with $\varphi\in\mathcal V$, $L:\IR\times\IR^{p}\to\IR^{p}$
Borel and $\varrho:\IR^{p}\to[0,1]$ Borel, write $\widetilde W(\varrho):=W-\varrho(\Delta Z)$ and define the
$\IR^{p}$-valued \emph{debiased moment function}
\begin{equation}\label{eq:psidef}
\psi(\theta,\zeta)\;:=\;\Delta Z\,\Delta U(\theta,\varphi)\;+\;\widetilde W(\varrho)\,
L\big(\Delta U(\theta,\varphi),\,\Delta Z\big),
\end{equation}
with $\Delta U(\theta,\varphi)$ as in \eqref{eq:trialres}. Write $\widetilde W:=\widetilde W(\varrho)$ at the
true propensity, $\zeta_\star:=(\vdagger,L_\star,\varrho)$ as in \cref{def:weight}, and
\begin{equation}\label{eq:psistar}
\psi_\star\;:=\;\psi(\theta_\dagger,\zeta_\star)\;=\;\Delta Z\,\Delta U+\widetilde W\,L_\star(\Delta U,\Delta Z).
\end{equation}
Define
\begin{equation}\label{eq:JSig}
J\;:=\;E\big[\Delta Z\,\Delta Z'\big]\;+\;E\big[\widetilde W\,\ell_\star(\Delta U,\Delta Z)\,\Delta Z'\big],
\qquad
\Sigma\;:=\;E\big[\psi_\star\psi_\star'\big] .
\end{equation}
\emph{The second summand carries $\ell_\star=\partial_uL_\star$ and not $L_\star$}, because $\theta$ enters
$\psi$ only through $\Delta U(\theta,\varphi)$ and differentiating $L_\star$ in that argument returns its
$u$-derivative. \emph{And under A1 that summand is the zero matrix, so $J=E[\Delta Z\Delta Z']$ exactly}
--- \cref{lem:momzero}\textup{(iv)}, whose whole content is
$E[\widetilde W\mid\Delta U,\Delta Z]=0$ applied to the $\sigma(\Delta U,\Delta Z)$-measurable factor
$\ell_\star(\Delta U,\Delta Z)\Delta Z'$. \eqref{eq:JSig} is nevertheless left in the two-summand form,
because $J$ is \emph{defined} as the $\theta$-derivative of \eqref{eq:psidef} --- which has two summands
whatever their values --- and because that is the form the estimator \eqref{eq:Jhat} takes, where the
second summand does \textbf{not} vanish.

\emph{$J$ must \textbf{not} be described as ``a different matrix'' from $E[\Delta Z\Delta Z']$}: wherever
\eqref{eq:JSig} is well formed the two are the same matrix, by \cref{lem:nusigma}\textup{(v)} at
$\Xi(u,z)=\ell_\star(u,z)\,z'$. \emph{What \textbf{is} true, and is why identification of
$\theta_\dagger$ must not be routed through \eqref{eq:JSig}, is the conditionality}: $\ell_\star$ exists
only under A2, so \eqref{eq:JSig} is a well-formed object only under A2, while identification of
$\theta_\dagger$ is A1b(iv)'s and holds whether or not A2 does.
% ARCHAEOLOGY.  A sentence asserting the opposite stood here through dispatch 12:
% "J is the matrix of A3 and it is not E[Delta Z Delta Z']".  The claim entered as
% a round-2 CORRECTION -- round 1 had routed identification of theta_dagger through
% A3, and the repair (M-3) routed it through A1b(iv) instead, on the ground that A3
% is "nonsingularity of a different matrix" and is itself conditional on A2.  The
% conditionality half of that ground is right and survives; the "different matrix"
% half is false, and it was asserted and never computed.  Its negation was already
% proved in this same file at lem:nusigma(v).

\emph{The typing note is about the \emph{form} and not the value.} $\ell_\star$ is $\IR^{p}$-valued and
$\Delta Z'$ is $1\times p$, so
$E[\widetilde W\ell_\star\Delta Z']$ is $p\times p$ and its \emph{form} is neither coordinatewise in the
$p$ coordinates of $\theta$ nor symmetric --- the typing note recorded at the head of Part~F, and what
makes the sum with $E[\Delta Z\Delta Z']$ well formed rather than a type error. Its \emph{value} is zero,
so $J$ itself is symmetric and positive semidefinite. \emph{The warning is therefore not idle; it
transfers to $\hat J$}, whose second summand in \eqref{eq:Jhat} is a sample average at
$\widehat{\widetilde W}^{(-k)}$ and $\hat\ell^{(-k)}$, does not vanish, and leaves $\hat J$
non-symmetric --- which is why \eqref{eq:Sighat} writes $\hat J^{-1}\hat\Sigma\hat J^{-1\prime}$ with the
transpose and not $\hat J^{-1}\hat\Sigma\hat J^{-1}$.
\end{definition}

% ===========================================================================
%  I.2  THE ESTIMATOR
% ===========================================================================

\begin{definition}[Cross-fitting, and the fixed anchor of the $\varphi$-step]\label{def:crossfit}
Fix $\mathsf F\ge2$, not depending on $n$, and a partition of $\{1,\dots,n\}$ into folds
$\mathcal I_1,\dots,\mathcal I_{\mathsf F}$ with $n_k:=|\mathcal I_k|$ and
$\lfloor n/\mathsf F\rfloor\le n_k\le\lceil n/\mathsf F\rceil$, chosen independently of the data. Write
$\mathcal I_{-k}:=\{1,\dots,n\}\setminus\mathcal I_k$, $O_i$ for the $i$-th observation vector of A0, and
\[
\IE_n[\varsigma]:=\frac1n\sum_{i=1}^{n}\varsigma(O_i),
\qquad
\IE_{n,k}[\varsigma]:=\frac1{n_k}\sum_{i\in\mathcal I_k}\varsigma(O_i) .
\]
\emph{At $\mathsf F=2$ this is Part~H's split, with $(\mathcal I_k,\mathcal I_{-k})=(\mathcal I_m,\mathcal I_a)$.}

\emph{Convention \textup{(P8)}, re-declared for Part~I --- and the parts of \textup{(P8)} do not have the
same standing.} \textup{(P8)}'s \emph{transfer rule} and
its two companions --- the $O_P$ companion and the $o_P$ companion --- are \cref{lem:pconv}(i)--(iii),
stated at a \emph{generic} $\sigma$-field, written $\mathcal G$ there and instantiated here, and they apply
with no extension.
% ARCHAEOLOGY.  Through dispatch 16 this clause read that the generic sigma-field
% was "unnamed there", which was true of the comment block those statements then
% lived in and is not true of the compiled lemma.
\textup{(P8)}'s \emph{conditioning convention} --- ``every expectation and every $O_P$ is
conditional on the auxiliary fold'' --- was stated \textbf{for Part~H}, at $\sigma(\mathcal I_a)$, and is
therefore \textbf{re-declared} here rather than inherited: throughout Part~I it is read at
$\sigma(\mathcal I_{-k})$, with $\mathcal I_k$ in place of $\mathcal I_m$. \emph{The licence is the same
either way, since the convention is a reading of a proof and not a hypothesis; the sentence is not, and this
definition is where the re-declaration lives.} \emph{That $\sigma$-field is written out at every site and is
given no symbol}: the
three natural candidates are all claimed --- $\mathcal G$ is \cref{lem:density}'s closed span,
$\mathcal E$ is the best-approximation error of \eqref{eq:jackson}, and $\mathcal F$ is the Fourier
transform of \eqref{eq:usobolev} --- and \cref{lem:balance} and \textup{(P8)} already write
$\sigma(\mathcal I_a)$ out, so the convention is inherited rather than invented.
Every conditional expectation and every conditional $O_P$ below is at that
$\sigma$-field, and A0 supplies the independence of $\mathcal I_k$ from $\mathcal I_{-k}$, which with
clause \textup{(c$'$)} below is what \textup{(P8)}'s \emph{freezing lemma} consumes.
\emph{Which clause of \textup{(P8)} each site uses is named at the site, and the answer for Part~I is
\textbf{never the transfer rule}.}
% ARCHAEOLOGY.  Through dispatch 13 this paragraph said otherwise.
Each conditional bound below is produced by conditional Markov or conditional Chebyshev, hence holds at
\emph{every} threshold with a $\sigma(\mathcal I_{-k})$-measurable numerator, and is converted by the
\emph{$O_P$ companion} when the numerator is unconditionally $O_P(1)$ after division by the stated rate
(\cref{lem:prelimtheta} Step~4, \cref{lem:sandwich} Step~2) and by the \emph{$o_P$ companion} when it is
unconditionally $o_P(1)$ (\cref{lem:epterm} Step~6, \cref{lem:sandwich} Step~8). \emph{The transfer rule
needs the conditional probability to vanish at a \textbf{single} moving threshold, which neither
inequality delivers}; that is the observation \textup{(P8)} already recorded for the $o_P$ case, now
applied to the $O_P$ case. What supplies the tightness of the numerators is A4(iii), A5(ii), (iv) and
hypothesis (R-a) of \cref{thm:rootn}.

For each $k$ the nuisances are computed on $\mathcal I_{-k}$:
\begin{enumerate}
\item[(a)] \emph{Propensity.} $\hat\varrho^{(-k)}$ satisfies A4(iii) with $n$ replaced by $n_{-k}$ and is
truncated to $[\varepsilon_\varrho,1-\varepsilon_\varrho]$ as in \cref{def:hatT}, so that
$\widehat{\widetilde W}^{(-k)}:=W-\hat\varrho^{(-k)}(\Delta Z)$ satisfies
$\varepsilon_\varrho\le|\widehat{\widetilde W}^{(-k)}|\le1$ surely and
$\mathsf r_k:=\hat\varrho^{(-k)}-\varrho$ satisfies $\|\mathsf r_k\|_\infty\le1-2\varepsilon_\varrho<1$.
\item[(b)] \emph{Structural function.} $\hat\varphi^{(-k)}$ is the estimator of \cref{def:lfstep}, run on
$\mathcal I_{-k}$ \textbf{at a fixed anchor $\theta_\circ\in\IR^{p}$ chosen before the data are seen}, and
$h_k:=\hat\varphi^{(-k)}-\vdagger$. \emph{The anchor is what removes the circularity, and it is licensed and
not free.} By \cref{cor:opsol} and \cref{lem:shiftinv} the population solution set of \eqref{eq:opeq} in
$\varphi$ is $\mathcal V_{CI}$ at \emph{every} $\theta$, so under A1b(iii) the $\varphi$-step targets
$\vdagger$ at $\theta_\circ$ as well as at $\theta_\dagger$; \emph{what is not $\theta$-free is the
conditioning of the inverse problem}, so \cref{prp:phirate}'s hypotheses (P-a)--(P-g), \eqref{eq:muhalf} and
A7 in the form \eqref{eq:a7amend} are read \emph{at $\theta_\circ$}, and $a$, $\mu$, $\mathsf R_\ast$ and
$\hat\delta_n$ are $\theta_\circ$-specific. This is \cref{rem:R1}\textup{(iv)}, used and not extended.
\item[(c)] \emph{Debiasing weight.} $\hat L^{(-k)}$ is an admissible weight in the sense of
\cref{def:weight}(b), $\hat L^{(-k)}=L^{\,\hat b^{(-k)}}_{\hat\ell^{(-k)}}$, generated by an estimator
$\hat\ell^{(-k)}$ of $\ell_\star$ that is Borel, bounded, and continuous in $u$ at every $(u,z)$;
$D_k:=\hat L^{(-k)}-L_\star$. \emph{Those three properties of $\hat\ell^{(-k)}$ are the minimum
\cref{def:weight}(b) needs to be well posed} --- they are what make $\int_0^u\hat\ell^{(-k)}(v,z)dv$ exist
and make $\partial_u\hat L^{(-k)}=\hat\ell^{(-k)}$ hold at every point --- \emph{and they are stated here
rather than inherited}, so that this definition is well posed on its own terms; hypothesis (R-a) of
\cref{thm:rootn} strengthens continuity to continuous differentiability, which \emph{one} clause of this
definition uses and only one --- \textup{(c$'$.2)} needs it to make the inner supremum of the second
envelope, $\sup_u\|\partial_u\hat\ell^{(-k)}(u,z)\|$, attained on a countable dense set and hence
measurable. \emph{So the sup-norm convention \eqref{eq:supnormconv} is \textbf{not} unconditional.}
% ARCHAEOLOGY.  Through dispatch 14 this sentence read "which no clause of this
% definition uses", from which a reader would conclude that it is unconditional.
\item[(c$'$)] \emph{Measurability of the five estimated nuisances, and the reading of their sup-norms.}
\emph{Clause \textup{(c)} above states a per-realisation ``Borel'' and nothing more; without the present
clause the hypotheses of (R-a) and the displays \eqref{eq:foldrates} are not well-posed probability
statements, because they are statements about the distribution of a supremum and of an integral of a random
function.} \emph{This clause re-declares nothing}:
\cref{lem:balance}'s sample-splitting convention --- the Part~F original --- states
$\sigma(\mathcal I_a)$-measurability of $\hat\ell$ as a \emph{random element} of
$L^{2}(P_{\Delta U,\Delta Z})$ together with a fixed \emph{jointly measurable representative}, which is
strictly weaker than \eqref{eq:prodmeas} below and is all that site needs, it evaluating $\hat\ell$ at an
\emph{undistorted} $(\Delta U,\Delta Z)$ draw.
% ARCHAEOLOGY.  Through dispatch 15 lem:balance's convention stated MEMBERSHIP and
% a measurability whose target was unspecified, and this sentence in turn described
% it as "a membership statement and not a measurability declaration".  Both were
% repaired in dispatch 16 at lem:balance.
\emph{And the right container is not the same for all five}, which is the substance of
\textup{(c$'$.1)}. Two clauses follow; neither is available from clause \textup{(c)} alone, and
\textup{(c$'$.1)} is a requirement on the estimators and not a consequence of anything above it ---
\emph{with \textbf{one} exception, marked at the clause itself, and it is the object clause \textup{(b)} has
already fully constructed}. For $\hat\varphi^{(-k)}$ the property is \textbf{verified} rather than required,
at the measurability paragraph of \cref{def:lfstep}; \emph{and the exception is forced, not stylistic}, since
one cannot require a property of an object one has already built --- \cref{rem:obligations}'s test, read on an
estimator.
% ARCHAEOLOGY.  Through dispatch 18 this sentence quantified over all five without
% the carve-out, so the clause simultaneously pinned hat-varphi^{(-k)} to a
% construction and demanded of it a property that construction was never shown to
% have.
\begin{enumerate}
\item[(c$'$.1)] \emph{Measurability --- and it is two different statements, because three of the five are
only ever integrated and two are evaluated pointwise at arguments their container does not control.}

\emph{The three integrated nuisances, as random elements of $L^{2}$ at a fixed representative.} Each of
$\hat\varrho^{(-k)}$, $\hat b^{(-k)}$ and $\hat\varphi^{(-k)}$ is a
\textbf{$\sigma(\mathcal I_{-k})$-measurable random element} of, respectively, $L^{2}(P_{\Delta Z})$,
$L^{2}(P_{\Delta Z})$ and $\mathcal V$ --- \emph{a measurable map from the sample space into the space, and
not a family of functions indexed by the sample point}. A random element of an $L^{2}$ space over a
separable metric space admits a jointly measurable representative --- measurable with respect to the
product of the sample $\sigma$-algebra with $\mathfrak B(\IR^{p})$ for the first two and with
$\mathfrak B(\mathcal X)$ for the third --- \emph{fix one for each, once, as A2c does for
$\ell^\Pi_\star$}.
\emph{For the third of the three this is a \textbf{discharged verification} and not a requirement, and the
difference is not cosmetic.} Clause \textup{(b)} pins $\hat\varphi^{(-k)}$ to \cref{def:lfstep}, whose
stopping index is \emph{random}, so the ground that settles the other two --- A4(iii)'s ``by inspection of the
formula'', one Part earlier --- is unavailable: no formula is inspected when an index is chosen by a
discrepancy rule. \emph{What supplies it instead is the measurability paragraph of \cref{def:lfstep}}, which
builds the property from the deterministic cap \eqref{eq:cap}, the emptiness convention of
\cref{def:lfstep}\textup{(c)}, and the finiteness of the resulting family: the stopping index is a hitting
index of finitely many measurable events, the iterates are finite compositions of measurable maps, and the
stopped iterate is a finite indicator-weighted sum of them. \emph{The superscripts are not written out here},
because \cref{def:lfstep}'s index is a \emph{stage} index and this clause's is a \emph{fold} index, and the
two are one glyph apart. \emph{The other two remain requirements, as printed}, and $\hat b^{(-k)}$'s is read
off its formula on the $\hat\varrho$ model.
\emph{The sample $\sigma$-algebra is written out and given no symbol}, because
calligraphic $\mathcal F$ is A2c's Fourier transform and $\mathcal F_{XZW}$ never appears
without its subscript; the same convention applies to $\sigma(\mathcal I_{-k})$
itself.
% ARCHAEOLOGY.  Glossary Sec.1.6, C15, and the C78 write-it-out precedent.
\emph{For these three the choice of representative is immaterial, and that is the test a container has to
pass --- but the test is passed on \textbf{two different grounds}.} For $\hat\varrho^{(-k)}$ and $\hat b^{(-k)}$ the evaluation is at
$\Delta Z_i$, whose law \emph{is} the one the container is built on, so identity of laws does it. \emph{For
$\hat\varphi^{(-k)}$ it is not}: its container is $\mathcal V\subset L^{2}(\mathcal X,\pi)$, whose null
sets are the $\pi$-null ones, while its evaluations are at $X_{i1}$ and $X_{i2}$, whose laws are $f_{X_1}$
and $f_{X_2}$. \emph{What transfers the null sets there is \textbf{domination}, and it has a tag}: C1(ii)
gives $\bar c_\pi=\max_{j\in\{1,2\}}\|f_{X_j}/\pi\|_\infty<\infty$, so every $\pi$-null set is
$P_{X_j}$-null for $j\in\{1,2\}$; \emph{``the law its container is built on'' is the wrong ground for the
third object}. Either way any two representatives
% ARCHAEOLOGY.  The domination ground was uncited through dispatch 15.
agree at all of the finitely many evaluation points outside a $P$-null set, the union of the finitely many
exceptional sets being $P$-null.

\emph{For $\hat\varrho^{(-k)}$ the choice is moreover \textbf{pinned}, and that is what clause
\textup{(a)}'s \emph{sure} bound rests on.} A4(iii), read at the fold by
clause \textup{(a)}, requires product measurability of $\hat\varrho$'s \emph{own pointwise function} ---
the \eqref{eq:prodmeas} shape, one Part earlier --- so for that object no equivalence class is formed and
no representative is chosen at all; \cref{def:hatT}'s truncation is then a pointwise operation on the very
function clause \textup{(a)} names, and
$\varepsilon_\varrho\le|\widehat{\widetilde W}^{(-k)}|\le1$ holds \emph{surely}, as clause \textup{(a)}
says. The $L^{2}$ container is retained
for $\hat\varrho^{(-k)}$ only because A4(iii)'s \emph{rate} is an $L^{2}$ statement; it is not what any
pointwise bound rests on.
% ARCHAEOLOGY.  Through dispatch 15, clause (a) asserted a SURE property of an
% object this clause had placed in an L^2 container, with nothing tying the chosen
% representative to the truncation -- this file's own test, applied to the wrong
% object.  Every consumer of clause (a) would have survived a weakening to "almost
% surely", so it was a wording defect and not a broken bound; it was repaired at
% A4(iii) in dispatch 16 because the assumption had to be repaired anyway.

\emph{The two weight objects, for which the $L^{2}$ container is not merely unnecessary but wrong.} What is
required of $\hat\ell^{(-k)}$, in addition to clause \textup{(c)}, is that its \emph{own} pointwise
function --- the one \textup{(c)} already defines, at every realisation and every $(u,z)$ --- be product
measurable:
\begin{equation}\label{eq:prodmeas}
(\omega,u,z)\;\longmapsto\;\hat\ell^{(-k)}(\omega)(u,z)
\quad\text{is}\quad
\sigma(\mathcal I_{-k})\otimes\mathfrak B\big(\IR\times\IR^{p}\big)\text{-measurable} ,
\end{equation}
and $\hat L^{(-k)}=L^{\,\hat b^{(-k)}}_{\hat\ell^{(-k)}}$ then \emph{inherits} the same property from
\eqref{eq:weightdef}, together with the representative fixed above for $\hat b^{(-k)}$ --- \emph{so
\eqref{eq:prodmeas} is imposed on one object and not on two}. \emph{The inheritance is the
\textbf{measurability half} of Fubini--Tonelli and it is taken on positive and negative parts,
coordinatewise}: \eqref{eq:weightdef}'s integrand is $\IR^{p}$-valued and
% ARCHAEOLOGY.  A dispatch-15 draft elided this to "by Tonelli's theorem applied to
% int_0^u hat-ell^{(-k)}(v,z) dv".
of arbitrary sign, and $\int_0^u=-\int_u^0$ when $u<0$. For each coordinate $j\le p$ and each sign part the
two maps
\[
(\omega,u,z)\;\longmapsto\;\int_\IR\un\{0\le v\le u\}\,\big(\hat\ell^{(-k)}_j\big)^{\pm}(\omega)(v,z)\,dv,
\qquad
(\omega,u,z)\;\longmapsto\;\int_\IR\un\{u\le v\le0\}\,\big(\hat\ell^{(-k)}_j\big)^{\pm}(\omega)(v,z)\,dv
\]
are product measurable by Tonelli's theorem, their integrands being non-negative and
$\sigma(\mathcal I_{-k})\otimes\mathfrak B(\IR\times\IR\times\IR^{p})$-measurable by \eqref{eq:prodmeas};
all four are \emph{finite} at every point because clause \textup{(c)} makes $\hat\ell^{(-k)}$ bounded at
each realisation, so their signed combination is well defined and product measurable. \emph{This is the
device A2c\textup{(m.2)} already owns}, one Part earlier and on the other side of the
pairing, where the same positive-and-negative-part split carries Tonelli across a signed
$u$-integral. \emph{Under (R-a) the derivative inherits it too, and \textup{(c$'$.2)} needs that for its
second envelope}: $\partial_u\hat\ell^{(-k)}$ exists at every point under (R-a) and is the pointwise limit
of the difference quotients of $\hat\ell^{(-k)}$ in $u$ taken along the reciprocals of the positive
integers --- \emph{the increment is written out and carries no letter}, on the same ground as
\eqref{eq:supnormconv}'s countable dense set; each such quotient is product measurable by
\eqref{eq:prodmeas}, and a pointwise limit of countably many product-measurable functions is product
measurable. \emph{A representative of an $L^{2}$ class supplies nothing of the kind}, having no derivative
in general.
\emph{No equivalence class is formed and no representative is chosen, so there is nothing
for a choice of representative to disturb, and clause \textup{(c)}'s continuity in $u$ is untouched.}
\emph{Why the $L^{2}$ reading fails here and not three lines above}: these two are evaluated at a
% ARCHAEOLOGY.  A draft of this clause declared the two weight objects random
% elements of L^2(P_{Delta U, Delta Z}) and asserted that "the continuity clause of
% (c) is preserved by the choice", which is false for an arbitrary representative,
% point evaluation not being continuous on L^2.
\emph{shifted} first argument --- at $(\Delta U-\lambda\Delta h_k,\Delta Z)$ in \cref{lem:biasexp} Steps~1
and~4--5, at $(\Delta U_i(\theta,\hat\varphi^{(-k)}),\Delta Z_i)$ in \eqref{eq:Jhat} and
\cref{lem:sandwich} Steps~2 and~9 --- whose law \textup{(c$'$.2)} below argues, and needs, \emph{not} to be
dominated by $P_{\Delta U,\Delta Z}$; on an undominated set two representatives of one class disagree, so
$\hat J$ of \eqref{eq:Jhat}, $\Psi^{(1)}_k$ of \cref{lem:biasexp}, and hence $\hat\theta$ of
\eqref{eq:onestep} and $\hat\Sigma$ of \eqref{eq:Sighat}, would depend on which was chosen and
\emph{the estimator would not be a function of the object the clause declares}.

\emph{What is retained from the $L^{2}$ reading}, because several results use it, is membership:
$\hat\ell^{(-k)}\in L^{2}(P_{\Delta U,\Delta Z})$ almost surely --- \cref{lem:balance} Step~0's third
point, and what makes the third display of \eqref{eq:foldrates} finite --- and, for $\hat L^{(-k)}$, the
same membership, which is \emph{derived and not required}: by \eqref{eq:weightdef},
$\|\hat L^{(-k)}(u,z)\|\le\|\hat b^{(-k)}(z)\|+\|\hat\ell^{(-k)}\|_\infty|u|$, so
$E\|\hat L^{(-k)}(\Delta U,\Delta Z)\|^{2}\le
2E\|\hat b^{(-k)}(\Delta Z)\|^{2}+2\|\hat\ell^{(-k)}\|^{2}_\infty E[\Delta U^{2}]<\infty$ almost surely
under \cref{def:weight}(b)'s admissibility, (R-a) and A1b(iv). \emph{Both statements are about the
function and neither is a container declaration.} \emph{The third display of \eqref{eq:foldrates} is a
statement with a truth value}: by \eqref{eq:prodmeas} and Tonelli's theorem,
$\omega\mapsto\|\hat\ell^{(-k)}-\ell_\star\|_{L^{2}(P)}$ is $\sigma(\mathcal I_{-k})$-measurable, and
likewise every conditional expectation and every $O_P$ statement below about these objects is well posed.
\emph{Equation \eqref{eq:prodmeas} is also what \textup{(P8)}'s freezing lemma consumes for the
$\hat\ell^{(-k)}$ \emph{ingredient} of its composite, at six of its seven Part~I sites}, that lemma being
stated at a composite of the form $(\omega,e)\mapsto f\big(\hat\ell^{(-k)}(\omega),e\big)$. \emph{The
seventh is \cref{lem:prelimtheta} Steps~4 and~4$'$, where $\hat\ell^{(-k)}$ does not appear at all and the
frozen object is $h_k$}: what supplies the composite there is the jointly measurable representative fixed
above for $\hat\varphi^{(-k)}$, so both halves of this clause are load-bearing and neither covers the
other. \emph{The composite is \textbf{not} product measurable ``precisely when'' \eqref{eq:prodmeas}
holds, in either direction.} Not sufficient: the composite is product
measurable only once $f$ is itself jointly Borel in its two arguments, and at the sites where the
evaluation point is built from $\hat\varphi^{(-k)}$ only together with that object's own clause. Not
necessary: $f\equiv0$ is product measurable whatever $\hat\ell^{(-k)}$ is. \emph{What is true is the
direction the proofs use}, and it is the one this clause is for: \eqref{eq:prodmeas} is what this
% ARCHAEOLOGY.  A dispatch-15 draft wrote the "precisely when".
definition \emph{supplies}, and it is what the per-realisation ``Borel'' of clause \textup{(c)} does not. \emph{The file owns the device by which \eqref{eq:prodmeas} is verified, and it
is recorded rather than imposed}: \cref{lem:biasexp} Step~1 obtains joint measurability in
$(\lambda,\omega)$ from ``continuous in $\lambda$, measurable in $\omega$'' by the Carath\'eodory
criterion, and the same criterion gives \eqref{eq:prodmeas} whenever, for each fixed $u$, the map
$(\omega,z)\mapsto\hat\ell^{(-k)}(\omega)(u,z)$ is
$\sigma(\mathcal I_{-k})\otimes\mathfrak B(\IR^{p})$-measurable, since \textup{(c)} already makes
$u\mapsto\hat\ell^{(-k)}(u,z)$ continuous at every $(\omega,z)$. \emph{That is a sufficient condition and
not the requirement}; the requirement is \eqref{eq:prodmeas}, which is what the proofs use.
\item[(c$'$.2)] \emph{The sup-norms are essential suprema in $z$ and genuine suprema in $u$.} For the two
envelopes (R-a) prices,
\begin{equation}\label{eq:supnormconv}
\big\|\hat\ell^{(-k)}\big\|_\infty
\;:=\;\operatorname{ess\,sup}_{z}\ \Big\{\ \sup_{u\in\IR}\ \big\|\hat\ell^{(-k)}(u,z)\big\|\ \Big\} ,
\qquad\text{the essential supremum taken under }P_{\Delta Z},
\end{equation}
and likewise for $\|\partial_u\hat\ell^{(-k)}\|_\infty$. \emph{Both halves of that reading are needed and
neither can be interchanged with the other.} The inner supremum is over an uncountable index but is
attained on a countable dense set, because \textup{(c)} makes $u\mapsto\hat\ell^{(-k)}(u,z)$ continuous at
every $z$ and (R-a) makes $u\mapsto\partial_u\hat\ell^{(-k)}(u,z)$ continuous as well, so
$(\omega,z)\mapsto\sup_u\|\cdot\|$ is jointly measurable by \textup{(c$'$.1)}.

\emph{The \textbf{outer} essential supremum is measurable as well, and this does \textbf{not} follow from
joint measurability without an argument --- it is not a countable supremum.} \emph{The argument is the
tool of the previous
paragraph, once more.} At each fixed real level, the map
% ARCHAEOLOGY.  A dispatch-15 draft of this clause carried the outer essential
% supremum on the word "likewise".
$\omega\mapsto P_{\Delta Z}\big\{z:\sup_u\|\hat\ell^{(-k)}(\omega)(u,z)\|>\text{that level}\big\}$ is
$\sigma(\mathcal I_{-k})$-measurable by Tonelli's theorem applied to the non-negative jointly measurable
indicator; the levels at which that probability vanishes are exactly the levels \emph{at or above} the
essential supremum, since it vanishes at the essential supremum itself --- the exceptional set there being
the union of the countably many exceptional sets one reciprocal-of-an-integer above it --- and at no
smaller level, by the definition of the essential supremum as an infimum; and therefore, for every real
level, the event that the essential supremum does not exceed it is the intersection, over the
\emph{countably many rational levels strictly above it}, of the events that the probability vanishes ---
a countable intersection of measurable sets, the \emph{strict} inequality being what makes those
exceptional sets exhaust the one at the level itself.
Hence $\omega\mapsto\|\hat\ell^{(-k)}\|_\infty$, and likewise
$\omega\mapsto\|\partial_u\hat\ell^{(-k)}\|_\infty$, is $\sigma(\mathcal I_{-k})$-measurable as a map into
$[0,+\infty]$. \emph{The levels are written out and carry no letter}, on the same ground as the countable
dense set above (C12/C24/C78). \emph{This is load-bearing and not hygiene}: (R-a) prices both envelopes as
$O_P(1)$, which is a statement about the distribution of these two maps, and \textup{(P8)}'s $O_P$
companion requires its \emph{bounding numerator} to be measurable with respect to the conditioning
$\sigma$-field --- and at \cref{lem:sandwich} Step~2 that numerator is
$\|\hat\ell^{(-k)}\|^{2}_\infty E\|\Delta Z\|^{2}$, whose only random factor is the envelope. \emph{No
symbol of the companion's is imported into compiled text}; the numerator is named by what it is.

\emph{A genuine supremum over $z$, by contrast, would not be a measurable function of $\omega$}, since no
clause makes $z\mapsto\hat\ell^{(-k)}(u,z)$ continuous. \emph{And the essential supremum must be taken in $z$ and not
in $(u,z)$}, because every consumer of \eqref{eq:supnormconv} evaluates the envelope at a
\emph{shifted} first argument --- at $(\Delta U-\lambda\Delta h_k,\Delta Z)$ in \cref{lem:biasexp} Steps~1
and~4--5, at $(\Delta U_i(\theta,\hat\varphi^{(-k)}),\Delta Z_i)$ in \cref{lem:sandwich} Steps~2 and~9 ---
whose law need not be dominated by $P_{\Delta U,\Delta Z}$, while its \emph{second} argument is in every
case a draw of $\Delta Z$. \emph{The exceptional set is therefore $P$-null where the envelope is used}:
for each $i$, $\sup_u\|\hat\ell^{(-k)}(u,\Delta Z_i)\|\le\|\hat\ell^{(-k)}\|_\infty$ almost surely, and
the union over the finitely many $i\le n$ of the exceptional sets is still $P$-null, so the bound holds
simultaneously at all $n$ observations almost surely. \emph{For the deterministic $\ell_\star$ no
convention is needed and none is imposed}: A$2'$(i)'s closing clause defines $\ell_\star$ pointwise and
\cref{cor:elldiff} bounds it at \emph{every} $(u,z)$, so $\|\ell_\star\|_\infty$ is a genuine supremum and
the two readings agree.
\end{enumerate}
\item[(d)] \emph{The fold reading of the maintained nuisance-rate clauses, declared once.} A4(iii) is read
at each fold in (a); the \textbf{four} clauses of A5 that price an \emph{estimated} nuisance are read the
same way, and every citation of them in Part~I is a citation of the corresponding display here:
\begin{equation}\label{eq:foldrates}
\max_{k\le\mathsf F}\big\|h_k\big\|_{L^{2}(\pi)}=O_P(\delta_n),
\qquad
\max_{k\le\mathsf F}\big\|h_k\big\|_{H^{s}}=O_P(1),
\qquad
\max_{k\le\mathsf F}\big\|\hat\ell^{(-k)}-\ell_\star\big\|_{L^{2}(P)}=O_P(\epsilon_n),
\end{equation}
being A5(ii)'s first half, A5(vi), and A5(iv)'s first half respectively --- \emph{the third in force precisely
when (R-e.1) of \cref{thm:rootn} is the branch taken, since under (R-e.2) no rate for $\hat\ell^{(-k)}$ is
assumed at all}. \emph{The second display is finite at every realisation only under
\eqref{eq:sieveinhs}, and it is cited here for the same reason A5(vi) carries it}: $h_k$ is
$\hat\varphi^{(-k)}-\vdagger$ with $\hat\varphi^{(-k)}$ an iterate of \cref{def:lfstep} and hence ---
\textbf{once $\hat\varphi^{(0)}$ is in $\mathcal S_n$} --- an
element of $\mathcal S_n$, so A8's basis clause is what places it in $H^{s}$ and H1 is what places
$\vdagger$ there; without the first the display is an $O_P(1)$ statement about a quantity equal to
$+\infty$ surely, which is a condition on the \emph{empty} event and not a rate.
\emph{The premise in bold is stated here because this is a \textbf{definition} and inherits no hypothesis
list, and because without it the sentence is false as written}: \cref{def:lfstep}\textup{(d)} sets
$\hat\varphi^{(k)}=\hat\varphi^{(k-1)}+v_{\hat m^{(k)}}$ with $v_m\in\mathcal S_n$, so the path lies in
$\mathcal S_n$ if and only if its \emph{starting point} does --- which is why A5(vi$'$), (P-a.1) and
\cref{def:lfstep}'s own opening all say ``once $\hat\varphi^{(0)}$ does''.
\emph{Every result that reads \eqref{eq:foldrates} supplies the premise},
(P-a.1) being in group (T-5) of \cref{thm:rootn}.
% ARCHAEOLOGY.  Through dispatch 20 this clause stated the display without the
% bold premise.
\emph{The indexing is not cosmetic
and it is not free either.} A5(ii), (iv) and (vi) are stated at a \textbf{single} estimated nuisance, so a
result that cites them bare while averaging over $\mathsf F$ fold complements is short by exactly the
maximum. For the first two the cost is nil: \cref{thm:rootn} Step~0 \emph{derives} them from (T-5) with
$\mathsf F$ fixed and $n_{-k}\asymp n$. \emph{For the third it is a strengthening of A5(iv) in form},
of the same order for fixed $\mathsf F$ but a strengthening nonetheless, and it is therefore carried
explicitly by hypothesis (R-e.1) of \cref{thm:rootn} rather than read into A5(iv). \emph{(R-e.2) and (R-a)
already write $\max_{k}$.}
\emph{The fourth clause is A5(vi$'$).} It prices the \emph{stagewise} Sobolev stability of the
$\varphi$-step,
% ARCHAEOLOGY.  (R-e.1) was left bare by a round-10 draft; this enumeration read
% "the three clauses" through dispatch 13.
$P(\max_{j\le\bar k}\|\hat\varphi^{(j)}-\vdagger\|_{H^{s}}\le\mathsf R)\to1$ --- where $\bar k$ is a
\emph{stage} index and not a fold index, the two being one glyph apart --- and it is a clause about an
\textbf{estimated} nuisance, so it is read at each fold complement exactly as the other three are: for each
$k\le\mathsf F$, A5(vi$'$)'s display is required of the $\varphi$-step run on $\mathcal I_{-k}$.
\emph{This costs nothing and is nonetheless not without consequence.} It costs nothing because $\mathsf F$
is fixed, so the fold reading is the intersection of $\mathsf F$ events each of probability tending to one;
it is not without consequence because the radius $\mathsf R$ must then serve all $\mathsf F$ fold
complements \emph{at one value}, $\mathsf R$ being a constant of \cref{def:srcball} and not a fold-indexed
object --- and $\mathsf R$ also appears in (P-c), (P-g), \eqref{eq:entropy} and, through
$\mathsf R_\ast$, in the schedule constants, so a fold-indexed radius is not available. (T-5) already reads
A5(vi$'$) at each fold \emph{operationally}, through ``\eqref{eq:phirate} holds on each
$\mathcal I_{-k}$''; what was missing is that this enumeration --- which declares itself the place where
the reading lives --- did not list it.
\end{enumerate}
Put $\hat\zeta^{(-k)}:=(\hat\varphi^{(-k)},\hat L^{(-k)},\hat\varrho^{(-k)})$ and, for $\theta\in\IR^{p}$,
\begin{equation}\label{eq:barpsi}
\bar\psi_n(\theta)\;:=\;\frac1n\sum_{k=1}^{\mathsf F}\sum_{i\in\mathcal I_k}\psi\big(\theta,\hat\zeta^{(-k)}\big)(O_i),
\qquad
\hat\psi^{(-k)}:=\psi\big(\theta_\dagger,\hat\zeta^{(-k)}\big).
\end{equation}
\emph{Since $\mathsf F$ is fixed, $n_k/n\to\mathsf F^{-1}$ and $n_k\asymp n$}; the constant $\mathsf F$ is
carried explicitly in every bound rather than absorbed, because \cref{prp:medianpart} quantifies over
partitions and a suppressed $\mathsf F$ would hide the one place a partition enters an order.
\end{definition}

\begin{definition}[The preliminary and the one-step estimator]\label{def:onestep}
With the notation of \cref{def:crossfit} define
\begin{equation}\label{eq:prelimtheta}
\hat\theta^{(0)}\;:=\;\Big(\IE_n\big[\Delta Z\Delta Z'\big]\Big)^{-1}
\frac1n\sum_{k=1}^{\mathsf F}\sum_{i\in\mathcal I_k}\Delta Z_i\big(\Delta Y_i-\Delta\hat\varphi^{(-k)}_i\big),
\end{equation}
on the event that $\IE_n[\Delta Z\Delta Z']$ is nonsingular, and $\hat\theta^{(0)}:=0$ otherwise. Define
the \emph{cross-fitted Jacobian}, for $\theta\in\IR^{p}$,
\begin{equation}\label{eq:Jhat}
\hat J(\theta)\;:=\;\frac1n\sum_{k=1}^{\mathsf F}\sum_{i\in\mathcal I_k}
\Big[\Delta Z_i\Delta Z_i'+\widehat{\widetilde W}^{(-k)}_i\,
\hat\ell^{(-k)}\big(\Delta U_i(\theta,\hat\varphi^{(-k)}),\Delta Z_i\big)\,\Delta Z_i'\Big],
\end{equation}
so that $\partial_\theta\bar\psi_n(\theta)=-\hat J(\theta)$ identically, and put
$\hat J:=\hat J(\hat\theta^{(0)})$. \emph{That identity uses \textbf{one} derivative and no more}:
$\partial_u\hat L^{(-k)}=\hat\ell^{(-k)}$, which \cref{def:weight}(b) supplies at every point under the three
properties \cref{def:crossfit}(c) requires of $\hat\ell^{(-k)}$, together with $\theta\mapsto\Delta U_i(\theta,\cdot)$
being affine by \eqref{eq:trialres} and the sum being finite. \emph{No second derivative is claimed here.}
$\partial_u^{2}\hat L^{(-k)}=\partial_u\hat\ell^{(-k)}$ is used only at \cref{thm:rootn} Step~1, for Taylor's
theorem in $\theta$, and it is hypothesised \emph{there}, at (R-a) --- \emph{a definition should not carry a
differentiability claim that a theorem assumes}. The \emph{one-step estimator} is
\begin{equation}\label{eq:onestep}
\hat\theta\;:=\;\hat\theta^{(0)}\;+\;\hat J^{-1}\,\bar\psi_n\big(\hat\theta^{(0)}\big)
\end{equation}
on the event $\{\hat J\ \text{nonsingular}\}$, and $\hat\theta:=\hat\theta^{(0)}$ otherwise. The variance
estimator is
\begin{equation}\label{eq:Sighat}
\hat\Sigma\;:=\;\frac1n\sum_{k=1}^{\mathsf F}\sum_{i\in\mathcal I_k}
\psi\big(\hat\theta,\hat\zeta^{(-k)}\big)(O_i)\;\psi\big(\hat\theta,\hat\zeta^{(-k)}\big)(O_i)' ,
\qquad
\hat V\;:=\;\hat J^{-1}\hat\Sigma\,\hat J^{-1\prime} ,
\end{equation}
the last \textbf{on the event $\{\hat J\ \text{nonsingular}\}$, with $\hat V:=0$ otherwise}.
\emph{The fallback is not optional.} $\hat\theta^{(0)}$ is given a value off $\{\IE_n[\Delta Z\Delta Z']\
\text{nonsingular}\}$ and $\hat\theta$ is given one off $\{\hat J\ \text{nonsingular}\}$;
$\hat V$ inverts the \emph{same} $\hat J$ and needs one too. That matters because
\cref{thm:rootn}'s \textbf{second} conclusion, $\hat V\xrightarrow{p}J^{-1}\Sigma J^{-1\prime}$, and the
studentised statement that follows it are statements about $\hat V$ at every $n$, and a convergence in
probability is not well posed for a quantity that is undefined on an event of positive probability.
\emph{The choice of value is immaterial and the zero matrix is chosen because it is the one value that
needs no further object}: by \cref{lem:sandwich}\textup{(i)}, $P(\hat J\ \text{nonsingular})\to1$, so the
event on which the fallback is active has probability tending to zero and no conclusion of
\cref{lem:sandwich}\textup{(ii)} or of \cref{thm:rootn} changes. \emph{In a finite sample the fallback is
not innocuous and should not be reported as a variance estimate}: $\hat V=0$ makes the studentised
statistic undefined, which is the honest report on that event and is the reason the value is stated rather
than left to the reader.

\emph{One step and not a root, and the ground is the \textbf{sample} criterion and not the population
one.} The realised map
$\theta\mapsto\bar\psi_n(\theta)$ is \emph{not} affine: by \eqref{eq:psidef} $\theta$ enters
$\hat L^{(-k)}$ through its first argument, and $\widehat{\widetilde W}^{(-k)}$ is not conditionally
centred, so $\bar\psi_n(\theta)=0$ need not have a unique root at any $n$. \emph{The ground is
\textbf{not} that no global rank condition is available}: by \cref{lem:momzero}\textup{(iv)} and
\eqref{eq:psiaffine} the
population moment $\theta\mapsto E[\psi(\theta,\zeta_\star)]$ is affine with gradient $-E[\Delta Z\Delta Z']$,
so a global rank condition \emph{is} available, it is A1b(iv), and $\theta_\dagger$ is its unique root on
all of $\IR^{p}$ --- which says nothing about the sample criterion at estimated nuisances.
% ARCHAEOLOGY.  Through dispatch 12 this paragraph gave the population half of the
% ground as "A3 is a LOCAL rank condition ... and no global rank condition is
% available".
\emph{What the block supplies is a preliminary with a
rate} (\cref{lem:prelimtheta}), and a one-step from a preliminary with a rate is exactly the construction
that converts a rank condition on the population Jacobian into a limit distribution without solving
anything; the cross-fitted form is that of
\citet{chernozhukov2018dml}, and the two classical anchors for the one-step itself --- Newey \& McFadden
(1994, Handbook ch.~36) \S3.4 and van der Vaart (1998) Ch.~5 --- are \emph{not cited here because their keys
are absent from this project's bibliography}; they are flagged for the librarian in
\texttt{theory\_memo.md} rather than written in, on the precedent \cref{rem:lfnonlinear}(c) sets for
Dunker (2021). \emph{The Jacobian is
evaluated at $\hat\theta^{(0)}$ and not at $\hat\theta$}, which is what makes \eqref{eq:onestep} a closed-form
estimator and what makes the algebra of \cref{thm:rootn} Step~1 exact rather than asymptotic.
\end{definition}

% ===========================================================================
%  I.4  THE THEOREM
% ===========================================================================

\begin{theorem}[\texttt{thm:rootn}: asymptotic normality of the cross-fitted one-step estimator]\label{thm:rootn}
The one-step estimator is $\sqrt n$-consistent and asymptotically normal with the sandwich variance, under
mild ill-posedness and a propensity smooth enough to put the structural rate strictly inside $n^{-1/4}$.
\emph{The hypothesis groups are tagged \textup{(T-1)}--\textup{(T-5)} and the additional hypotheses
\textup{(R-a)}--\textup{(R-e)}, on the model of \cref{prp:phirate}'s \textup{(P-a)}--\textup{(P-g)}.}
% ARCHAEOLOGY.  A draft of this theorem tagged the groups (H1)--(H5).  H1 is the
% Hilbert-scale link, it is cited INSIDE group (T-5), and a group label reading
% "H1" next to an assumption tag reading "H1" is the collision this file has
% repaired four times elsewhere.
Suppose
\begin{enumerate}
\item[(T-1)] \emph{(sampling, model, identification)} A0; A1; A1b(i), A1b(iii), A1b(iv);
\item[(T-2)] \emph{(the representer and the weight)} A2; A$2'$(i.a), A$2'$(i.b), A$2'$(iii); N1 at
$L_\star$. \emph{A3 is \textbf{not} in this group}, being proved from
(T-1) and the rest of (T-2) at \cref{lem:momzero}\textup{(iv)}; the rank condition the theorem needs is
A1b(iv), which is in (T-1);
% ARCHAEOLOGY.  A3 stood in this group through dispatch 12.
\item[(T-3)] \emph{(index set, densities, propensity)} B1(i), (ii), (iv); C1(ii); C2$'$(i); A4(i),
A4(ii), A4(iii);
\item[(T-4)] \emph{(rates --- and the group is \textbf{not} homogeneous in status, which is why it is
annotated)} A5(i) and A5(ii) in the form \eqref{eq:deltalog} below, \emph{both of which (T-5)
\textbf{delivers} rather than requires} --- A5(i) through \cref{lem:opnoise} and A5(ii)'s first half through
\eqref{eq:phirate}, per Step~0 of the proof --- listed here because Steps~2 and~4 cite them by name and
because \eqref{eq:deltalog} fixes the object the symbol denotes; A5(iii), \emph{which no result in this file
delivers and which is therefore assumed}, being a restriction on the primitives $(\mu,r,d_Z)$ that
\cref{prp:phirate} presupposes rather than establishes --- that proposition returns \eqref{eq:phirate}
\emph{whenever} \eqref{eq:frontier} holds and says nothing when it fails; and A5(vi) with $s>\tfrac12$, read
at each fold as in \cref{def:crossfit}(d), of which (T-5)'s A5(vi$'$) delivers the $H^{s}$ bound at the
\emph{final linearisation stage} $\bar k$ of \cref{def:lfstep} --- \emph{a stage index, not a fold index},
the two being one glyph apart in this very clause --- while $s>\tfrac12$ is a restriction on the primitive
$s$;
\item[(T-5)] \emph{(the $\varphi$-step)} the hypotheses of \cref{prp:phirate}, read at the anchor
$\theta_\circ$ of \cref{def:crossfit}(b) --- H1; A5(vi$'$); A8's basis and upper-end clauses;
$p_K\ge s-\tfrac12$; (P-a.1)--(P-a.3), (P-c)--(P-g); and \eqref{eq:muhalf}, $\mu>\tfrac56$ strictly ---
so that \eqref{eq:phirate} holds on each $\mathcal I_{-k}$;
\end{enumerate}
and suppose in addition
\begin{enumerate}
\item[(R-a)] \emph{(the weight is built as an anchored antiderivative, with a bounded derivative)} for each
$k$, $\hat L^{(-k)}=L^{\,\hat b^{(-k)}}_{\hat\ell^{(-k)}}$ in the sense of \cref{def:weight}(b), with
$\hat\ell^{(-k)}$ continuously differentiable in $u$ at every point and bounded,
\[
\max_{k\le\mathsf F}\Big\{\big\|\hat\ell^{(-k)}\big\|_\infty+\big\|\partial_u\hat\ell^{(-k)}\big\|_\infty\Big\}=O_P(1),
\qquad
\max_{k\le\mathsf F}E\big\|\hat b^{(-k)}(\Delta Z)\big\|^{4}=O_P(1),
\]
and $\max_{k\le\mathsf F}\|\hat L^{(-k)}-L_\star\|_{L^{2}(P)}=o_P(1)$.
\emph{\textbf{The \emph{sure} boundedness of $\partial_u\hat\ell^{(-k)}$ is \emph{not} required}}: no
first-step construction the block contemplates delivers a sure bound on an estimated object, what a
series or kernel estimator of $\ell_\star$ supplies being a bound that is finite with probability tending
to one --- exactly the $O_P(1)$ display above. The two sites that read finiteness of the derivative
envelope, \cref{lem:biasexp} Step~4 and \cref{lem:sandwich} Step~5, carry it on the \textbf{right} of an
inequality and are read on the event where it is finite, which is where every $O_P$ conclusion in those
two lemmas is read in any case. \emph{Boundedness of
$\hat\ell^{(-k)}$ itself is retained and is not part of this}: it is what \cref{def:crossfit}\textup{(c)}
already requires, for \cref{def:weight}(b) to be well posed. \emph{And \cref{def:crossfit}\textup{(c$'$.2)}'s $[0,+\infty]$ typing is unchanged by
this}: that typing is the \emph{definition}'s, fixing the range of the map $\omega\mapsto\|\cdot\|_\infty$
% ARCHAEOLOGY.  Dispatch 19 added a sure-boundedness conjunct on the ground that
% the two proof steps read finiteness straight off this hypothesis, and the same
% paragraph then recorded that "nothing was broken by the omission".  Both cannot
% be true of a minimal clause; the conjunct went and the O_P(1) display stayed, and
% the two proof sentences were re-grounded on the display rather than deleted.  A
% third site, thm:rootn Step 1, reads the clause only for the C^2 property of
% hat-L^{(-k)} and for the O_P rate, and is untouched either way.
before any hypothesis is imposed, and it must stay extended-valued because \cref{def:crossfit} is adopted by
results that do not list (R-a); \emph{tightness} of that extended-valued map --- and hence finiteness of it
with probability tending to one, but \textbf{not} at every sample point --- is what this clause adds to it,
at the results that do;
\item[(R-b)] \emph{(fourth moments, and exactly two of them)} $E\|\Delta Z\|^{4}<\infty$ and
$E[\Delta U^{4}]<\infty$. \emph{A fourth moment of $b_\star$ is \textbf{not} needed}: $b_\star$ enters
Part~I at two lines only, both at \emph{second} moment through
A$2'$(iii) (\cref{lem:momzero} Step~1), and nothing in Part~I needs a fourth moment of $\psi_\star$ ---
\cref{lem:sandwich} Step~8 applies conditional Markov to squares, its Step~10 needs only
$\IE_n\|\psi_\star\|^{2}=O_P(1)$, and the central limit theorem used is Lindeberg--L\'evy. The
\emph{fourth}-moment need is on $\hat b^{(-k)}$, and (R-a) supplies it;
% ARCHAEOLOGY.  A round-10 draft carried a third clause, E||b_star||^4 < infinity.
\item[(R-c)] \emph{(the $\varrho\times L$ cross term)} \textbf{either} \textup{(c.1)} $D_k=\hat L^{(-k)}-L_\star$
satisfies N1$'$ for every $k$, \textbf{or} \textup{(c.2)}
$\max_k\|D_k\|_{L^{2}(P)}=O_P(\epsilon^{L}_n)$ with $r_n\,\epsilon^{L}_n=o(n^{-1/2})$;
\item[(R-d)] \textbf{\sout{SKIPPED.}} There is \textbf{no} hypothesis (R-d): the list runs (R-a), (R-b),
(R-c), (R-e). The
letter is \emph{reserved and not reused}, on the \textup{(P-b)} precedent, so that a citation to
``(R-d)'' lands on this line and not on the wrong hypothesis;
\item[(R-e)] \emph{(the first-order bias)} \textbf{either} \textup{(e.1)} A5(iv) read at each fold, i.e.\ the
third display of \eqref{eq:foldrates} together with $\delta_n\epsilon_n=o_P(n^{-1/2})$ --- \emph{the fold
index is part of the hypothesis and not a reading of it}, for the reason \cref{def:crossfit}(d) gives, and it
is what \cref{thm:rootn} Step~4(a) and \cref{lem:sandwich} Step~3 consume --- \textbf{or}
\textup{(e.2)} the hypotheses of \cref{lem:balance}\textup{(iii)} --- the sieve solve sets term
\textup{(II)} of \eqref{eq:balance3} to zero, A8 holds with a B-spline basis of degree exceeding
$s-\tfrac12$, A5(vi) holds and \eqref{eq:nuhat} holds --- together with
$\max_k\|\hat\ell^{(-k)}-\ell_\star\|_{L^{2}(P)}=o_P(1)$, \textbf{and} with
\begin{equation}\label{eq:re2window}
s\;>\;\frac{1}{2\,(1-2\alpha)} .
\end{equation}
\emph{\eqref{eq:re2window} is stated on this branch because this
branch is the one place both of its ingredients arrive at once.} \Cref{lem:balance}\textup{(iii)} carries A8 \emph{in full}, so (e.2) imports A8's lower
end $L_x\gg n^{1/(2s)}$; (T-5) simultaneously imports (P-d), $L_x=o\big(n^{1-2\alpha}/\ln n\big)$; and
\eqref{eq:frontier} forces $\alpha>\tfrac14$, so (P-d) binds and the joint window for $L_x$ is non-empty if
and only if \eqref{eq:re2window}. \emph{The arithmetic is printed at (P-d) already} --- ``had A8's lower end
been carried as well, the window would instead require $s>(2(1-2\alpha))^{-1}$'' --- and
\cref{rem:h6bandwidth}\textup{(5)} records that the reference design clears it. \emph{(T-4)'s only
$s$-floor is $s>\tfrac12$, which does not reach it.}
\emph{Witness, chosen so that nothing else is empty there}: $(\mu,s,a,\alpha)=(2,1.3,0.325,0.32)$ clears
\eqref{eq:frontier} at $\kappa\alpha=0.256$, clears \eqref{eq:lxwindow} at $s>1$ and clears
\eqref{eq:hxwindow} at $s>1.2273$, while A8's lower end asks $L_x\gg n^{0.3846}$ and (P-d) allows
$L_x=o(n^{0.36}/\ln n)$: \emph{empty}, the threshold \eqref{eq:re2window} being $s>1.3889$.
\emph{So without \eqref{eq:re2window} the branch is vacuous below that threshold and the theorem is
\textbf{not}}: (e.1) is unaffected and carries the conclusion there. \emph{An either/or that advertises a
choice which does not exist at part of the parameter set is a defect of the enumeration, not of the
theorem}, and \eqref{eq:re2window} is what makes the choice real wherever it is offered.
\end{enumerate}
Then, with $\hat\theta$ of \eqref{eq:onestep} and $J,\Sigma$ of \eqref{eq:JSig},
\begin{equation}\label{eq:rootn}
\begin{split}
\sqrt n\,\big(\hat\theta-\theta_\dagger\big)
&\;=\;J^{-1}\,\frac1{\sqrt n}\sum_{i=1}^{n}\psi_\star(O_i)\;+\;o_P(1)\\[2pt]
&\;\xrightarrow{\;d\;}\;\text{the centred Gaussian law on }\IR^{p}\text{ with covariance matrix }
J^{-1}\Sigma J^{-1\prime} ,
\end{split}
\end{equation}
and $\hat V\xrightarrow{p}J^{-1}\Sigma J^{-1\prime}$, so that for \textbf{every} $\vartheta\in\IR^{p}$ with
$\|\vartheta\|=1$,
\[
\sqrt n\,\vartheta'\big(\hat\theta-\theta_\dagger\big)\big/\big(\vartheta'\hat V\vartheta\big)^{1/2}
\;\xrightarrow{\;d\;}\;\text{the standard Gaussian law on }\IR .
\]
\emph{The quantifier is ``every'' and not ``every $\vartheta$ with
$\vartheta'J^{-1}\Sigma J^{-1\prime}\vartheta>0$''.} By
\cref{lem:momzero}\textup{(v)} and \eqref{eq:sigmapd}, $\Sigma\succ0$ under A1b(iv) and C2$'$(i) --- both
already in (T-1) and (T-3) --- so $J^{-1}\Sigma J^{-1\prime}\succ0$ and the positivity guard is
\emph{implied by the theorem's own hypotheses}. \emph{A proviso that a theorem's hypotheses entail should
be stated as a consequence and not as a condition the user checks.} \emph{What the guard is \textbf{not} replaced by is a
% ARCHAEOLOGY.  Through dispatch 13 the quantifier read "every vartheta with
% vartheta' J^{-1} Sigma J^{-1'} vartheta > 0"; the change is a strengthening of
% the statement at no cost in hypotheses, and the entailment was found by probing
% this very display.
uniform lower bound}: $\lambda_{\min}(\Sigma)>0$ at the fixed $P$ of the theorem and may be arbitrarily
small at nearby laws, which is \cref{rem:rootnscope}\textup{(6)}. \emph{The denominator $\vartheta'\hat V\vartheta$
can still be zero or negative in a finite sample} --- $\hat V$ is a sample sandwich and is set to $0$ off
$\{\hat J\ \text{nonsingular}\}$ by \eqref{eq:Sighat} --- so the display is read, as such displays always
are, on the event of probability tending to one where the denominator is strictly positive.
\emph{The contraction direction is $\vartheta$, which is this file's binding letter for exactly this
object}: $\vartheta'\psi$ is the moment for $\vartheta'\theta_\dagger$. It is \textbf{not} $\varkappa$,
which is a local bound variable in \eqref{eq:jackson}, nor $\mathsf r$, which is the
$\varrho$-direction.
% ARCHAEOLOGY.  Glossary C25.
\emph{Both limit laws above are written out and carry no symbol}: italic $N$ is a count, a null set or a
covering number everywhere else in this corpus, and calligraphic $\mathcal N$ is
unavailable, being \cref{lem:density}'s closed subspace and A2c's null Borel set.
% ARCHAEOLOGY (dispatch-23 change to how this theorem is PRINTED and to nothing it
% says).  Through dispatch 22 the two displays read N(0, J^{-1} Sigma J^{-1'}) and
% N(0,1) -- a law with two arguments on the letter lem:pfreeze's clause reserves.
% Step 5 of the proof already said "the limit law is written out and carries no
% symbol" while this statement printed the symbol twice.  The census behind glossary
% C84 named its own pattern as "N(", which does not match "N\big(", so the two
% displays here and the third at prp:medianpart were invisible to it: the proof was
% repaired and the three STATEMENTS were not.  Glossary C90 prints the per-file
% counts.
\end{theorem}

\begin{corollary}[T2$'$: the rate floor, without \eqref{eq:frontier}]\label{cor:ratefloor}
The same decomposition delivers a rate whether or not the frontier holds; \emph{on the sieve branch, and
only there, every exponent in that rate is a function of the primitives}.
Suppose (T-1)--(T-3), (T-5), (R-a), (R-b) hold, together with A5(i), A5(vi) with $s>\tfrac12$,
$\delta_n\to0$, $\epsilon_n\to0$, $r_n\to0$, and (R-c.2), (R-e.1) with the two products left
\emph{unrestricted}. Then
\begin{equation}\label{eq:ratefloor}
\big\|\hat\theta-\theta_\dagger\big\|
\;=\;O_P\Big(n^{-1/2}\;+\;\delta_n\epsilon_n\;+\;\delta_n^{2}\;+\;r_n\delta_n\;+\;r_n\epsilon^{L}_n\Big).
\end{equation}
Under (R-e.2) in place of (R-e.1) the term $\delta_n\epsilon_n$ is replaced by
$O_P(n^{-1/2}\delta_n)+O_P(\mathsf R\,L_x^{-s})$, and no rate for $\hat\ell$ enters
\eqref{eq:ratefloor} at all.

\emph{How much of \eqref{eq:ratefloor} is in primitives, stated by branch, because the answer differs by
branch.}
% ARCHAEOLOGY.  A round-10 draft of this corollary claimed the stronger of the two
% answers on the weaker branch.
\begin{enumerate}
\item[(a)] \emph{Always in primitives.} $\delta_n=(\hat\delta_n\ln n)^{2\mu/(2\mu+1)}$ with
$\hat\delta_n=n^{-r/(2r+d_Z)}$ and $r_n=n^{-r/(2r+d_Z)}$, so the terms $n^{-1/2}$, $\delta_n^{2}$ and
$r_n\delta_n$ have exponents that are functions of $(\mu,r,d_Z)$ alone --- through \cref{prp:phirate} and
A4(iii).
\item[(b)] \emph{Never in primitives on the branch stated above.} $\epsilon_n$ is priced by A5(iv) and
$\epsilon^{L}_n$ by (R-c.2), and \emph{neither has a supplier anywhere in the maintained block}: $\hat\ell$
is constructed nowhere (\cref{rem:rootnscope}\textup{(3)}) and $\epsilon^{L}_n$ is priced by nothing at all.
So on the hypotheses as listed --- (R-e.1) with (R-c.2) --- the terms $\delta_n\epsilon_n$ and
$r_n\epsilon^{L}_n$ carry two exponents that are \textbf{assumed and not derived}, and the claim ``the rate
is in primitives'' is false for them. \emph{Note also that $s$ enters no exponent on this branch}, A8 being
consumed only through (R-e.2).
\item[(c)] \emph{Entirely in primitives on one branch, and that branch is what \S19.1's sieve decision
buys.} Under (R-e.2) \textbf{and} (R-c.1) --- the sieve solve for the first-order bias, N1$'$ at each
$D_k$ for the $\varrho\times L$ term --- the two unsupplied symbols both disappear, and
\[
\big\|\hat\theta-\theta_\dagger\big\|
=O_P\Big(n^{-1/2}+n^{-1/2}\delta_n+\mathsf R\,L_x^{-s}+\delta_n^{2}+r_n\delta_n\Big),
\]
every exponent of which is a function of $(\mu,r,d_Z,s)$ through \cref{prp:phirate}, A4(iii), A8's window
and the spline direct theorem. \emph{That is the sense in which this is a result and not bookkeeping, and it
holds on that branch only.}
\end{enumerate}
\emph{No further index enters.} A2c's $\mu''$ supplies a finite \emph{constant} through
\cref{lem:reprsmooth} and \cref{cor:elldiff} and appears in no exponent, and the project's second
smoothness index is $\mu''$ and no other.
% ARCHAEOLOGY.  A round-10 draft of the sentence above listed a symbol "mu-prime",
% defined nowhere in this project and with no glossary row; it was deleted rather
% than defined.
Finally, \eqref{eq:rootn} is the special case of \eqref{eq:ratefloor} in which every term after the first is
$o(n^{-1/2})$.
\end{corollary}
